
%%%%%%%%%%%%%%%%%%%%%%%%%%%%%%%% SCRIPT FOR E-RARA AND MDZ FILES     %%%%%%%%%%%%%%%%%%%%%%%%%%%%%%%%%%%%%%%%%%%%%%%%
%%%%%%%%%%%%%%%%%%%%%%%%%%% fini le 30.04.2025 par F. GOY            %%%%%%%%%%%%%%%%%%%%%%%%%%%%%%%%%%%%%%%%%%%%%%%%
% !TeX TS-program = lualatex
\documentclass{article}
\usepackage[T1]{fontenc}
\usepackage{microtype}
\usepackage[pdfusetitle,hidelinks]{hyperref}

\usepackage{polyglossia}
\setmainlanguage{english}
\setotherlanguages{latin,greek}
\usepackage[series={},nocritical,noend,noeledsec,nofamiliar,noledgroup]{reledmac}
\usepackage{reledpar}

\usepackage{fontspec}
\setmainfont{TeX Gyre Termes}

\usepackage{sectsty}
\usepackage{xcolor}

\usepackage{fancyhdr}
\pagestyle{fancy}
\fancyhf{}
\fancyhead[LE,RO]{\nouppercase{\leftmark}}  
\cfoot{\thepage}
\renewcommand{\headrulewidth}{0.4pt}

% Redefine \section to remove numbering
\usepackage{titlesec}
\titleformat{\section}[block]{\normalfont\scshape\color{gray}}{}{0pt}{} % no number in heading
\titleformat{\subsection}[hang]{\normalfont}{}{0pt}{} % also remove subsection number
\titleformat{\subsubsection}[hang]{\normalfont\footnotesize\color{black}}{}{0pt}{}

% Modify how section marks are stored to exclude numbers
\makeatletter
\renewcommand{\sectionmark}[1]{%
	\markboth{#1}{}} % Only store the section title, without number
\renewcommand{\subsectionmark}[1]{%
	\markright{#1}} % Only store the subsection title, without number
\renewcommand{\numberline}[1]{} % Hide the section number in TOC
\makeatother

\begin{document}

\date{}
        \title{Commentarii in epistolas d. Pauli: [Lefèvre d'Etaples, Jacques ], [1512]}
\maketitle
\tableofcontents
\clearpage
\begin{pages} 
\beginnumbering
        
\section*{COM. }
\marginpar{[ p.]}
\phantomsection
\addcontentsline{toc}{subsection}{\textit{¶COMMENTARIORVM IN EPISTOLAS BEATI PAVLI APOSTOLI LIBER DECIMVS QVI EST IN EPISTO. PRI. AD TIMOTHEVM. CAP. I. }}
\subsection*{\textit{¶COMMENTARIORVM IN EPISTOLAS BEATI PAVLI APOSTOLI LIBER DECIMVS QVI EST IN EPISTO. PRI. AD TIMOTHEVM. CAP. I. }}
\section{CAP. I.}\pstart \huge\textbf{P}\normalsize RIMA haec ad Timotheum epistola priuata est: sed publicam continet institutionem. Et praemittit vt solet sanctam salutationem apostolus dicens. ¶Paulus apostolus IHESV CHRIS TII secumdum ĩperium dei saluatoris nostri et domini IHESV CHRISTI spei nostrae: Timotheo ingenuo filio in fide gratia /misericordia/ pax a deo patre et CHRISTO IHESV domino nostro. ¶ Hac sancta salutatione posita: causam Timotheo aperit cur eum relinquere voluerit Ephesi cum transmisso Aegeo mari profectus est in Macedoniam. et narrationem: litera satis continet aptam. Dicit itaque . ¶Sic rogaui te vt remaneres Ephesi cum proficiscerer in Macedoniam: vt praeciperes quibusdami ne varie docerenti neq intenderent fabulis/ et genealogiis interminatis/ quae questiones pariunt magis q̃ gubernationem dei in fide. ¶ Quid est ne varie docerent: nisii ne aliter quam euangelicum dogma quod a spiritu sancto est docerent ? Praeterea iudaei multas fabellas habebant/ et genealogias indeterminatas: in quibus volebant potius intendere eos quos instituebant quam ad diuinae legis intentionem atque finem. Non est enim finis legis:scire huiusinodi figmenta qualia sunt innumera in Talmuth iudaeorum neque  genealogias eorum quae a nullo ad amussim sunt traditaei sed altercationes generant non aedificationem. Etiau si genealogiam Sem vsque ad Abraham ex genesil et eandem ex Luca deducere velis: non nisi altercationem pariet et nodum quem ex traditis dissoluere non possis, cuius rei ex septuaginta: concordia cognoscitur. et caput geneseos vbi describitur genealogia Sem mancum essel atque ita nobis esse legendum. Vixitque Sem postquam ge nuit Arphaxat. 5oo amnis et genuit filios et filias et mortuus em. Porro Arphaxat vixit annis 135 et genuit Cainan. vixitque Arphaxat postquam genuit Cainant amnis 430 et genuit fi lios et filias et mortuus em. Cainan cum 13o esset amnorum genuit Sale. Et reliqua. Haec inq̃ finis et praecipua intentio non sunt legis: sed perfecta dei et proximi dilectio. vt finis medi cinae est sanitas: non occupatio circa fabulas vanas Chironis. Esculapii Machaonis Po3 dalirii et eorum genealogias. Et qu i medicina em sanitas:i lege est dilectio. ꝗ enim perfectam » habet dilectionem: non aliter illi lege opus estiq̃ sano medicina. Ideo subdit Paulus, ⁋Finis autem praecepti est dilectio ex corde puroi et conscientia bona et fide non ficta. A quibus quidam aberrantes: ad vaniloquium conuersi sunt.ꝗ volunt esse legis doctores: cum non intelligant. .neque q̃ loquuntur neque de quibus affirmant. ¶ Quae dilectio perfecta sit et quam lex intendit: ostendit, est enim dilectio puri cordis conscientiae bonae et fidei non fictae, si enim cor purum non esset: quo modo esset perfecta dilectioquae est cordis et aiae spiritualis sanitas quandoquidem in corpore ĩpuroi neque est perfecta corporalis sanitas? Conscientia bona: mens est conscia boni et quae solum bonum et vere bonum vult. vt ex opposito conscientia mala: mens est con scia malil et quae malum cupit, non cum hac conscientia: sed cum illa est perfecta dilectio. Si etiam simulas te fidem habere vt perfidus Torquatus simulabat se esse Christianum: vt fidelis Tiburtius qui eum in fide illuminare voluerat raperetur ad tormenta aut idolis immolaret si inquam simulas fidem: eum simulata fide non est perfecta dilectio. sine enim fide: impossibile est placere deo. et cum simulata fides non sit fides: vir simulatae fideil deo non placet immo displicet. non igitur dilectionem habet: quae est munus dei quod ipse elargitur amicis. Vt igitur perfecta dilectio in nobis esse possit: sit cor purum sit conscientia bona sit fides syncaera et non simulata, vt cor purum et mumdum habeamus: a deo efflagitemus qui summa est puritas. vt conscientiam bonam: ab eodem qui summa est bonitas. fidem etiam non simulatam ab eodem: quia summa est veritas. Qui cor purum non habent conscientiam bonam( et fidem syncaeram aberrant et vana loquuntur: vt imperiti medicit nescientes et quae loquuntur et de quibus affirmant. Quo in genere maxime delinquunt qui iudaizant. loquuntur enim de lege: sed non intelligunt. nam spiritualis est: ipsi vero carnaliter concipiunt. et figuris inhaerent: de eis affirmantes quae debentur exemplaribus, tribuunt vmbrae: quae debentur veritati. Ergo dices: lex mala est? Nequaq mala est: nisi male vtentibus  et male intelligentibus. vt et medicina bona est: sed bene intelligentibus atque bene  \pend
\section*{X }
\section*{1 Timo. }
\marginpar{[ p.202 ]}\pstart ,vtentibus. Quod et subdit hoc pacto Paulus. (Scimus autem quod bona est lex:si quis ea ,legitime vtatur, qui sciat q iusto lex posita non est: sed iniquis et rebellibus impiis et ,peccatoribus scelestis et prophanis patris et matris percussoribus) homicidis fornica,toribus cinaedis plagiariis mendacibus periuris et si quid aliud sanae aduersatur do,ctrinae. quae est secundum euangelium gloriae beati dei: cui ego creditus sum. ¶Sic facit Paulus: vt qui diceret medicinam non esse positam pro sanis sed pro aegris et variarum egritudinum corporis enumeraret species. sic et varios aiae morbos enumerat:aduersus quos lex po sita fuiti vt variis medicamentis purificationibus et sacrificiis variis occurreret morbis  si i eos legis transgredientes praeceptalinciderent, ea enim lex vetabat: quia animae noxia. vt regulae medicorum vetant quae sunt isalubria. Qui iniqui id est anomi dicantur: notum est. sunt eni diuinae legis transgressores. Qui rebelles. nam sunt qui obedire nolunt i hoc ordinationi diuinae obluctantes. Impii: quibus deus non aliter curae est atque si non esset. Peccatores nomen generale est: etsi maxime manifestos et publicanos compraehendit. Patris et matris percussores noti sunt. et vtinam sic non essent: vt noti essent. sub quibus patri et matri con tumeliosi continentur: qui patrem et matrem non honorant quin potius irreuerentur et conuiciis  iniuriisque lacessunt. Homicidae plusquam par est et deceret noti sunt: ꝗ et vtinam nusq̃ noti essent  non verentes diuinam violare imaginem et quoad possunt penitus delere. Fornicatores dia boli sues: qui carnem immunde diuidunt alienae carni adhaerentes. et hi: fornicarii carnis. Alii sunt: qui spiritum diuidunt qui alieno spiritui q̃ dei adhaerent. et hi:fornicarii spiritus. Auaro; spiritui auaritiae adhaeret. Intemperans: spiritui luxuriae et gulael qui imundi sunt spiritus. Ambitiosus honorum: adhaeret spiritui superbiae. et omnes hi fornicatores spiritus: quatenus diuidunt spiritum. Cinaedi: in naturam peccant et deum qui author est naturae. et adeo infames sunt:vt natura nomina eorum exhorrescat neque ferme eos nominet nisi per circumlocutionem, quae et malet numq̃ nomina inuenta fuisse. nunq̃ enim inuenta essent: nisi natura (serpente antiquo obtinente) viam suam corrupisset. Plagiarii sunt: qui hoies praedantur vt redemptioni adigant aut i mancipia vendant. quales pyratae circa ripas maris: vt circa Massiliam Montempessulanum Anconant Apuliam esse feruntur et sunt. Mendaces et periuri hac pessima ifamiae nota no ti sunt. mendacia et periuria: falsa sunt numismatal quibus summae veritatis offenditur ma iestas, quae ideo aeterno igni et aeternis tormentorum bullitioniby cum suis authoribus tradentur: nisi antea purgentur q̃ in iudicium offensi regis veniant. his oibus et inumeris aliis quae sanctimoniae euangelicae aduersantur: varia ex lege prouisa erant remedial et pene inumera. At pro omnibus illis: Christus solus est vnica vera et sufficientissima medicinai si i . agens quod ad euangelii officium commissus siti atque dicit. ¶ Et gratias ago ĩualescere me fa,cienti CHRISTO IHESV domino nostro: quod fidelem me duxit ponendo in ministe,rio qui prius eram blasphemus et persecutori et iniurius. Sed misericordiam consecu,tus sum: quia ignorans feci in infidelitate. supermultiplicata est autem gratia domini ,nostri cum fide et dilectione: quae est in CHRISTO IHESV. Fidelis sermo et omni ,acceptatione dignus:quod CHRISTVS IHESVS venit in mundum vt peccatores sal,uos faceret. quorum primus: ego sum. Sed ideo misericordiam consecutus: vt in me primo osten,deret IHESVS CHRISTVS omnem longanimitatem ad informationem eorum qui ad vitam aeternam credituri sunt in eum. ¶ Ecce Paulus cum fuerat blasphemus persecutori iniurius infidelis et id deot licet ignorans: cum venit ad cognitionem suorum delictorumi non confugit ad remedia legis ad victimas et caerimonias sed ad solum Christum qui omnium peccatorum est vera medicina. Et blasphemorum in Christum persecutorum ecclesiael et nominis eius iniuriorum et contumeliosorum primus fuit: qui cognoscitur tam aperta misericordia vocatus. ideo inter omnes blasphemos persecutores Christi et contumeliosos: se primum vocat peccatorem qui misericordiam dei consecutus est. quod ideo a dei bonitate factum dicit: vt faceret eum caeteris peccatoriby exemplum longae expectationis suae erga eos(vt eos ad misericordiam suscipiat et eis suam elargiat gra tiam et suae dilectionis feruorem. Quam imensam et admirabilem erga nos Christi et dei beni8 , ¶¶ Regi autem saeculorum immortali inuisibili soli sapienti deo: honor et gloria in sae ,cula saeculorum. Amen. ¶Nullus rex terrenus id potuit praestare hominibus nullus mor talis nulla visibilis creatura snulla sapientia potuit inuenire creata:sed solus ille qui rex saeculorum omnium est per se immortalis inuisibilis et per se sapiens. Hac laude et gratiarum actione posita: hortatur Timotheum vt secundum praecedens ei indultum do2  tus et veritatis fidei ex spiritu prophetie ei cognita atque ait. ¶Hoc praeceptum commen. eum credimus si eum diligimus et in eo confidimus. Q uod et Paulus aperit: pariter gratias gnitatem Paulus consyderans: erumpit in tantae bonitatis laudem et gratiarum actionem atque subdit. num prophetiarum: viriliter militiam Christi exerceat declarando et docendo ea quae sunt spiri\pend
\section*{COM. }
\marginpar{[ p.1 ]}\pstart do tibi fili Timothee: vt secundum praecedentes in te prophetias milites in eis bonam ,militiam fidem habens et bonam conscientiam, quam quidam abiicientes: in fide naufra gium fecerunt. , ex quibus est Hymenaeus et Alexander: quos tradidi sathanaei vt discant non blasphema,re. ¶ A domino accepit Timotheus tam praeceptum quod ei commendat Paulusq̃ prophetiarum donum. Donum prophetiae duplex. Qddam est ad praedicendum futura: ad aliquorum salutem et aedifi cationem. Qu dam autem: ad itelligendans scripturas prophetarum. Itaque prophetiae quas habuit Ti motheus et q̃ magis aedificationi erant necessariae: erat intelligentia praeteritorum prophetarum de Christo et de sancta ad omnium gentium salutem in eum fide. hae enim prophetiae: fidei erant. Quid nam aliud factum putamus apostolis: cum dominus aperuit eis sensum vt intelligerent scripturas nisi spiritum prophetiarum eis infusum per quem omnes intelligerent prophetas? de Hy menaeo autem et Alexandro quod eos tradidit sathanae: Chrysostomus excommunicationem et resecationem a corpore Christi intelligit. ita vt nulla cum eis fidelium esset communio nulla spiritualium suffragiorum participatio. Et haec est citra interitum et vexationem carnis et occulta sathanae traditio. Illa autem quae fuit Corinthii incestuosi: manifesta fuit, cum interitu et vexatione carnis terrifica. Et id non vt perirent: sed vt ex hac deiectione a sanctorum societate discerent non blasphemarel non male loqui de Christol et eius fidel et resipiscentes dignamque poenitentiam agentes iterum insererentur corpori Christi. a quo ꝗ separatus esti vt ouis a grege sepata et praesente lupol ilico em sub lupi potestate: ita et ille continuo em sub sathane (qui diabolus et serpens antiquus est) potestate. A cuius morsibus solus Christus eripit: qui et nos semper protegat defendat eripiati rex saeculorum in omnia saecula benedictus. Amem.  \pend\pstart  ꝗEXAMINATIO nonnullorum circa literam. ¶ Vulgata aeditio. Timotheo dilecto filio in fidel gratia et misericordia et pax. (LVocabulum quod hic vetus interpres vertit dilecto: ad Philippenses cap.4 vertit germane. vbi dixit, etiam rogo et te germane conpar designatque proprie ingenuuml proprium ac peculiarem etiam authore Hieronymo. Et haec vocabula gratia miseriçordia pax:absque coniunctione in Paulo leguntur. Paulus. τιμοθέωο τμησίῳ τέκρω ἐμ Gίςςι, xαpις, ἐλςοσ, ἐιρη̂pη. ¶Vulgata aeditio.  Sicut rogaui te vt remaneres Ephesi. (Potius sic q̃ sicut: legendum. Neque enim similitu dinem comparationemue facit sed potius intendit: sic hoc pactol tali lege te rogaui quo remaneres Ephesse vt praeciperes quibusdam ne varias iducerent doctrinas et reliqua quae apostolus expressit. Paulus. κατὼς ωαρεκάλεσάσε ωροσμεί͂φαι ἐμ ἐφέσῷ. ¶¶ Vul . gata aeditio. Ne aliter docerent. ¶ Potius: varie. Paulus. μὴὶ ἐτερολιλασκαλεῖν. , ¶ Vulgata aeditio. Quae questiones praestant magis q̃ aedificationem deilquae est in fide. ¶Potius dicendum:q̃ dispensationensiue gubernationem. Paulus. αιτιμεσ τητησεις παρέχουσι μαλλομ, ἠι ὀικορομίαμ θεοῦͅ τὴμ ἐμ τίςει. Quid enim haec oeconomia est: nisi dispensatiol regiment et gubernatio domus deil in rectae fidei alimonial ¶ Vulgata aeditio. Si qs ea legitime vtatur scientes hoc quod lex iusto non est posita. ¶ Dicendum est singulariter sciens: et reuocatur ad quis non ad nos. quapropter ne allucinatio fieret, per relatiuum qui et verbum accommodate pro locutione resoluimus sciens ꝗ sciat. Paulus. ἐάν τισ ἀυτῶ ρομίμως χρῆται, ἐἰλὼσ τοῦτο, ὀτι λικαίσͅ pόuoσ oυ Rεῖται. Et quid mirum si iusto lex posita non est: quandoquidem saluator ait ? Non est , opus valentiby medico: sed male habentibus. ¶ Vulgata aeditio. Sceleratis et contaminatis) patricidis et matricidis homicidis fornicariis masculorum concubitoribo, plagiariis  , mendacibus et periuris. ¶ Potius dicendum: patris et matris percussoribus. Nam de percussione patris et matris: lex manifeste in Exodo expressa est. Qui percusserit (iquit) pa trem suum aut matrem: morte moriatur. De patricidio et matricidiolexpressa lex non esti quasi id tam enormes tam a natura abhorrendum sit: vt nunq contingere debeat. Et de omni genere flagiciosorum quos Paulus nominat: illic leges expressas comperies si diligenter exquirere volueris. Et plagiarii fures surreptoresque virorum sunt: sic dicti quod eis pedes ligent et in vincula coniiciant dum eos in mancipia vendant, aut forte se redimant. quales contingit plurimos esse eiusmodi hominum fures tempore bellorum. qd̃ quidem vocabulum graecuml q̃ apud nos latinum magis expromit. Et sine coniunctione dicitur mendacibus periuris. Paulus. αpοσ ίοισ, καὶ ἀφεφή λoις, πατραλοίαισ, καὶ μητραλοίαισ, ἀμAροφόpoισ. πόρpoισ, ἀρσε̃͂μokοίταις, ἀμAρατολις αῖσ, dέυςαις, ἐττι όpkοισ. ¶ Vulgara ,aeditio. Quod creditum est michi. (¶Sermo conuersus, nam legitur: cui creditus sum ego. Sed in idem recidit. Paulus. o ἐτιθε υθμμ ἐτω. ¶LVulgata aeditio. Gratias ago er qui me confortauit in CHRISTO IHESV domino nostro. ¶Dicendum est. Et gratias agolicet graece habeatur: et gratiam habeo. Nam consuetudo latina id habet: cum de deo loquimur. Et id confortauit: designat vires praestititi inualescere fecit. Et quod dicit in Christo Ihesu domino nostro: fine praepositione dici debet et eiusdem casus  \pend
\section*{1 Timo. }
\section*{x }
\marginpar{[ p.203 ]}\pstart esse cum ei. est enim is cui Paulus gratias agit:et qui eum inualescere fecit qui et ei vires for titudinemque suggessit. Paulus. καί χάριν ἔχω τῶ ἐνδυμώσαντί με χριστῶ ἰησοῦ  κυρίω ἡμῶν. ¶LVulgata aeditio. Sed misericordiam dei consecutus sum. ¶ Dei: abundati  etsi nichil intelligentiae officit. Paulus. ἀλλὰ ἠλεήθην. ¶ Vulgata aeditio. Venit in  hunc mundum. ¶Satis est dicere:in mundum. Paulus. ἦλθεν εἰς τὸν κόσμον. ¶ Vulgata aeditio. Ostenderet CHRISTVS IHESVS. ¶Mutatus vocabulorum ordo. dicit enim  Ihesus Christus. Paulus. ἐνδείξηται ἰησοῦς χριστός . ¶ Vulgata aeditio. Regi autem  saeculorunt imortalit inuisibilit soli deo. ¶ Dicendum: soli sapienti deo. Deest enim sapienti. Paulus. τῶ δὲ βασιλεῖ τῶν αἰώνων, ἀφθάρτω, ἀοράτω, μόνῶ θεῶ. Vt solus sapiens: quia solus nichil permixtum ignorantiae habens. sic solus rex saeculorum: quia solus omnium dominatur. solus immortalis: quia solus absolute interire non potest caetera autem possunt. solus inuisibilis: quia caetera compraehendi possunt (compraehendi enim est quodam modo videri) ipse vero solus compraehendi non potest. Et vt sacer ait Cyprianus. Videri non potest: visu clarior est. nec compraehendi: tactu purior est. nec aestimari: sensu maior est. Et ideo sic eum digne aestimamus: dum inaestimabilem ducimus. Ergo cum Iacob cum Moses dicumtur deum vidisse: id non designat deum appraehendisse sed aliquod vestigium et certum dei indicium attigisse. non eam ipsam incomprae hensibilem quae incompraehensibilis omnia appraehendit dei naturam atque substantiam.  \pend
\endnumbering\beginnumbering\section{CAP. II.}\pstart \huge\textbf{N}\normalsize Vnc quo modo ecclesiae ad publicam pietatem sint instituendae: informat Timotheum apostolus ac ait. ¶ Obsecro igitur primum omnium: vt fiant preces orationes supplicationes gratiarum actiones pro omnibus hoibus pro regibus et omnibus qui in excellentia sunti vt tranquillam et tacitam vitam degamus in omni pietate et honestate. Nam hoc bonum et acceptabile in conspectu saluatoris nostri dei: qui omnes homines vult saluos fieril et ad agnitionem venire veritatis. ¶Ante omnia vult publicam pietatem praecedere: et id deo gratum esse et acceptabile qui vult omnes homines saluos fieri et ad agnitionem venire veritatis. et si venirent: omnes hoies in magna pace et tranquillitate vitam degere possent. si venirent inquam: effectu ipso ad cognitionem veritatis. Nam sunt qui verbo fatentur et facto veritatem negamr: et illi nondum ad perfectam cognitionem sufficientemque venerunt ad tranquillam et tacitam degendam vitam ineque in omni pietate et honestate viuunt. Sunt alii qui neque verbo neque facto veritatem cognoscunt: veritatem dico quae necessaria est ad salutem. et haec non est dei voluntas: sed propriae voluntatis eorum a veritatis luce incuruatio. Nam qui vnum non cognoscunt deum et vnum Christum dominm nostrum pro omni salute mortem oppetiisse non fatentur et mediatorem inter deum et homines vt qui medium sit per quod omnes homines quotquot saluabuntur saluari oportet: ii nondum ad agnitionem venerunt veritatis etsi omnem scientiam et philosophorum et magorum et gymnosophistarum haberent exploratam. Haec enim omnia scire: si ad illud conferas nichil est scire. haec omnia sciens: ad immensas ignorantiae deuoluitur tenebras efficiturque tandem perpetuus ignorans. illud vero sciens: ad imensum sapientiae lumem extollitur efficiturque omnia sciens et sciens quidem cum vitae et sapientiae thesauris, cuius diuitiarum nullus est finis. illi vero: cum ignorantia cumulus erit indeficiens omnis in morte paupertatis. Hanc itaque diuitem sapientiam omni mundi sapientia maiorem: his paucis verbis explicat Paulus dicens. ¶ Vnus enim est deus. enim et mediator dei et hominum homo CHRISTVS IHESVS: qui dedit semetipsum redemptionem pro omnibus. ad qd̃ testimonium propriis  temporibus positus sum ego praeco et apostolus (veritatem dico in CHRISTO haud memtior) doctor gentium in fide et veritate. ¶ Ad testimonium illius ditissimae sapientiae q̃ vitam et salutem praebet animabus in mansionibus lucis aeternae in plenitudine temporum a deo definita: positus est Paulus praedicator et legatus non a seipso (nam nullus apostolus nullus legatus: a seipso) sed ab eo cuius est legatus. itaque a deo et Christo: a quo et doctor gentium factus esti non ad vllam disciplinam mundil quae nichil facit ad salutent sed ad doctrinam coelestem  ad doctrinam spiritus quae sola saluat et quam angeli magnifaciunt et docerent si quam docerent. doctor (inquit) gentium in fide et veritate. Quid hoc in fide et veritate: nisi in fide vera in fide veritatis in fide filii dei qui est via veritas et vita? Non igitur se effert non gloriatur cum dicit veritatem dico haud mentior: positus sum doctor gentium in fide et veritate  sed testimonium reddit veritati ei gloriam dans qui fecit eum praeconem suum apostolum  et doctorem gentium. Cum itaque sit constitutus a deo doctor gentium et doctoris officium sit docere : quod nunc eum viros? tum mulieres doceat audiamus. Docet itaque primum viros hocpacto   dicens. ⁋ Volo igitur vt viri in omni loco orent: leuantes sanctas manus sine ira et contentione. ⁋Quid est in omni loco: nisi in omni ecclesia vbicumque locorum in terra cognoscitur Christus? Sine ira et contentionc. Nam et dominus mandat: si offers munus tuum ad altare et ibi  \pend
\section*{COM. }
\marginpar{[ p.2 ]}\pstart recordatus fueris quia frater tuus habet aliquid aduersum tetrelinque munus tuum ante altare et vade prius reconciliari fratri tuo. et tunc veniens: offeres munus tuum. Quo modo enim audes te praesentare domino pacis: non habens signum pacis in mente tua sed oppositum signum quod est discordiae? Ergo nullo irato concitatoque animos neque qui dissentionem cum altero habet oratione maxime publica orare audeat. priuata autem vt deus adiuuet iraque sedetur et dissentio cesset pacemque suam restituat: orare nichil prohibet. is enim facit: quod aeger ad medicum suum languorem detegens et opem implorans. et haec:non publica oratio est quae intercessio dicitur quam sine ira et contentione deo offerre oportet. sed priuata: et quaedam (si sic dicere mauis) confessio suique accusatio. Docet et mulieres dicens. ¶ Identidem et mulieres in amictu mundolcum verecundia et modestia ornantes seipsas (non i tortis crinibus aut aurotaut margaritis aut veste sumptuosa: sed quae  decet mulieres profitentes diuinam pietatem per bona opera. Mulier in silentio discat:  cum omni subiectione. Mulieri autem docere non concedo:neque viri authoritatem habere   sed esse in silentio. Nam Adam primus plasmatus est: deinde Eua. Et Adam non seductus  est: sed mulier seducta in praeuaricatione fuit. Saluabitur autem per filiorum generationens  si manserint in fide et dilectione et sanctitate cum modestia. ¶ Haec sancti apostoli doctrina pro mulieribus. Mulierum habitus mundus sit:non sumptuosus. ornamentum mulierum sit verecundia et modestiat siue castitas. talis enim mulieres spirituales decet ornatus: adiunctis bonis operibus per quae diuinam pietatem profiteri videantur. Cunctis mulieribus ex inobedientia primae mulieris ingenita esse debet verecundia quae eas retrahat a perpetratione malorum vetitorum. ad quae ex imitatione primae parentis secundum carnem pronae feruntur. Et ex virgine culpae parentis Euae reparatrice: debet omnibus mulieribus virginitatis amor et appetitus inesse. quam si cum foecunditate habere non possunt:studere tamen debent habere foecundam castitatem quae virginitatis proxime aemula estErgo ex prima natiuitate verecundiam: ex secunda virginitatem aut saltem castitatem habere debent. hanc: theosebias id est erga deum pietatis et purus ac syncaerus dei cultus sequitur. illam vero: cautio fugaque malorum. Mulier non putet ornamentum in crinibus tortis aut aurosaut margaritis aut veste sumptuosa:qualis luxus nostra tempestate multas occupat cum alibil tum in vrbibus insubriae si forte contingit illuc iter facere. quae manifeste filiae Belial potius iudicari possunt: quam filiae Christil nisi quis oculos praestigiatos. et fascinatos habeat. Mulier si ad virum conferatur qui caput est: haec vt appetitus estrille vero vt ratio. at appetitui non conceditur vt doceat: sed rationi. ita neque mulieri: sed virorett maxime in publico. appetitus: a ratione formatur et discit. ita et mulier: a viro discere debet. mulier dico: quae vxor sit. quamquam et omnes in publico orationis loco:a viro qui communem omnium habet informationis vicem audire possunt in silentio. non quod temerarie veniant etiam virgunculae ad os publici praeceptoris: aliqua quesiturae. interrogent enim aut patres aut matres aut sapientiores matronas: quae melius quam ipsae potuerunt publicam vocem intelligere. Mulier: nunquam ausit authoritatem supra virum caperet id est dominari proprio viro. Nam hoc diuinae ordinationi reluctatur atque voluntati eius qui eam posuit sub viri potestate et virum fecit dominum, sub viri (inquit) potestate eris: et ipse dominabitur tui. O generis insotentia et opulentia dotis: quot mulieres fecit stultas quae authoritatem virorum sibi vendicanti seipsasi sexumque suum et dei voluntatem ignorantes superbae stultae insanae quantuncumque sublimes se existiment non hominibus sed deo rebelles? qui enim diuinae ordinationi resistit: deo resistit. estque deo rebellis: similis angelo tenebrarum qui hac de causa e coelis ruit in abyssum. Et mulierem viro deferre debere: vel ex ipsa creatione constat. Nam vir primo creatus esttanquam primogeniturae sibi vendicans honorem. mulier vero secundo loco. Et mulier seducta fuit: non vir. persuasa enim fuit a serpente: non vir. absente viro: deserpsit fructum vetitum, mandit: et viro eius rei ignaro porrexit mandendum. mandit porrectum: vnde generi nostro misera irrepsit mortis contagio. Attamem cum magis obnoxia culpae q̃ vir teneret: voluit nichilominus deus tam viro q̃ mulieri lalutem esse communem. mulieri quidem ob liberorum procreationem (ad id enim necessaria viro est) et id maxime dum vir et mulier in fide manserint et dilectione etsanctitate cum modestia atque castitate. castum enim tandiu thorum habent: quamdiu carnem suam per fornicationem neque hicineque illa diuidunt. Alioqui si fides non assit non dilectiot non sanctimonial vitaeque puritas et thori castitas: procreatio filiorum mulierem non saluabit. Nam preciosa ad saluandum pignora fides dilectio sanctitas castitas: magis quam carnalis proles. Id tribuente domino nostro omnium saluatore: qui erat et est et erit benedictus sine fine et in omnia saeculorum saecula. AMEN.  \pend\pstart  ¶ EXAMINATIO nonnullorum circa literam. ⁋ Vulgata aeditio. In omni sanctitate et  castitate. ⁋ Vocabulum quod hic vetus interpres castitatem dicit:et castimoniam et grauitate  \pend
\section*{1 Timo. }
\section*{x }
\marginpar{[ p.204 ]}\pstart et honestatem pro loco interpretari solent. Paulus.ἐν πᾶσι ἔυσεβεία, καὶ σεμνότητι. ¶ Vulgata aeditio. Cuius testimonium temporibus suis confirmatum est: in quo  positus sum ego. ⁋ Potius dicendum est: ad quod testimonium. Et abundat confirmatum est. Paulus. το μαρτύριον καίροις ἰδίοις ἐις ὃ ἐτέθην ἐγω. illud enim ἐις ὃ  reuocaui ad orationis caput. ¶ Vulgata aeditio. Veritatem dico non mentior. ¶Codices Graeci habent: veritatem dico in Christo. Paulus. αλήθειαν λέγω ἐν χριστῶ , ὀυ ψέυδομαι. ¶Vulgata aeditio. Leuantes puras manus sine ira et disceptatione, ⁋ Sunt qui volunt hic positum vocabulum dubitationem et haesitationem significare. verum si a διαλέγομαι venit: non male disceptationem disputationem dissentionem et contentionem quandam significat. Et cum iram et contentionem inuicem habemus: quod tunc manus ad deum non sint leuandae praeceptum est saluatoris. Si offers (inquit) munus tuum ad altare et ibi recordatus fueris quia frater tuus habet aliquid aduersum te relinque ibi munus tuum ante altare et vade prius reconciliari fratri tuo: et veniens offeres munus tuum. Cum enim puros conspectui diuino nos praesentamus: exaudit nos pro clementia sua qui audire exaudire et superexaudire vult vt par est et oportet paratos. Paulus. ἐπάιροντας ὁσίους χεῖρας, χωρὶς ὀργῆς καὶ διαλογισμοῦ. ⁋ Vulgata aeditio. Cum verecundia et sobrietate ornantes se. ⁋ Quod hic vetus interpres dicit sobrietate: vsitatius modestia dici potest. Vult etiam sacer Hieronymus castitatem dici ac modestiam. Paulus. μετὰ ἀιδοῦς καὶ σωφροσύνης κοσμεῖν ἐαυτάσ. ¶ Vulgata aeditio. Sed quod decet mulieres promittentes pietatem per bona opera. ¶ Dicendum esset quae foeminine:non quod neutro genere. Nam hoc relatiuum refert veste. quod ex graeco cognoscitur vbi masculinum est referens ἱματισμῶ masculinum. Et non generale pietatis nomen ponitur sed speciale: quod solum ad deum desi gnat pietatem. Paulus. ἀλλ ὅς πρέπει γυυαιξὶν ἐπαγγελομέναις θεοσεβειαν  Δἰ ἔργων ἀγαθῶν. ¶Vulgata aeditio. Neque dominari in virum: sed esse in silentio. ¶Potius:neque viri siue in virum authoritatem habere. Paulus. ὀυδὲ ἀυθεντεῖν ανδρός, ἀλλ’ εἶναι ἐν ἠσυχία. Verum authoritatem viri siue in virum habere et dominari viro:idem possunt. At quot sunt nunc matronae ex genere tumentes aut ex diuitiis opulentioreque dote quae nunc in viros authoritatem habenti facientes dei sententiam irritam propter fortunae ventum ventoque vaniorem fortunam. O caeca fortuna: quae non solum caeca sed quae vniuersum caecauit saeculum et cum qua veritas apostolica lucere  non potest. ¶Vulgata aeditio. Si permanserit in fide et dilectione et sanctificatione cum  sobrietate. ¶ Et hic quoque dici potest cum modestia aut castitate: vt in prima huius capitis examinatione dictum est. Et in principio particulae dicendum est pluratiue: manserint siue permanserint, respicit enim communiter vtrunque Paulus et virum et mulierem. Paulus. ἐὰν μείρωσιν ἐν πίστει, καὶ ἀγάπη, καὶ ἀγιασμῶ, μετὰ ςωφροσυνης.  \pend
\endnumbering\beginnumbering\section{CAP. III.}\pstart \huge\textbf{V}\normalsize T doctor gentium docuit apostolus viros et mulieres qui funt parentes carnales:nunc episcopos docet qui parentes sunt spirituales. et quales esse debeant et quae eorum officia ostendit, dicuntur enim episcopi a superintendendi officio. superintendendi dico in omni spiritualitate vitae: quod sanctum opus atque officium est. Ac proinde dicit. ¶ Fidelis sermo. Si quis episcopen desyderat: bonum opus desyderat. Oportet igitur episcopum: irrepraehensibilem esse vnius vxoris virum vigilem castum ornatum hospitalem docilem  non vinosum non percussorem non turpis lucri sectatorem: sed mitem imbellem non auarum propriae domui bene praesidentem: filios habentem in subiectione cum  omni honestate (si quis autem propriae domui praeesse nescit: quo modo ecclesiae dei curam administrabit?) non nouicium: ne superbia elatus in iudicium incidat diaboli. Oportet autem ipsum etiam testimonium bonum habere apud eos qui foris sunt:ne in vituperationem incidat et laqueum diaboli. ⁋ Episcope: visitatio intendentia obseruatio episcopi in suos dicitur. sicut episcopus: explorator intendens obseruator et custoai vt pastor gregis aut vt in alta specula vigilans custos ciuitatis. pulchrum sane nomem:sed rei nominis executio pulchrior quod officium episcopi est. quod cum tam pulchrum et venerandum sit: episcopum quam illi officio illique operi praeficitur irrepraehensibilem esse oportet, repraehensio: per omnia vicia currit superbiam iram ĩuidiam acediam gulam luxurians auaritiam. si aliquo horum viciorum notatus est aut aliqua specierum quae innumerabiles sunt: non est irrepraehensibilis. Quare neque  omnino dignus: huic committi officio. Bigamus qui duas duxisset vxores:non est idoneus  vt assumatur in episcopum. Nam carnem suam diuisit: et vnitatem in symbolo non custodiuit q̃ vnitatem Christi et ecclesie representat. Ideo in illa vnitate: sancti fiatuerunt eum non praeficiendum  \pend
\section*{COM. }
\marginpar{[ p.3 ]}\pstart qui in imagine violauit vnitatem id est qui illius spiritualis vnitatis violauit simulacrum. Q si bigamus qui in matrimonio et casto thoro vnitatem diuisit repellitur et censetur hac presi dentia indignus: quid igitur de iis ꝗ in fornicationibus carnem suam diuiserunt non solum bifornicatores sed qui centum (miserabile dictu) forte diuersarum pellicum fornicationiby se polluerunt toties carnem suam diuidentes et carnis suae vnitatem miserabiliter scidentes? Si illi repellum tur: et hi magis repellendi sunt. Argumentum apud Paulum efficax et apud omnes: qui esse norunt a minori ad maius. At canones hos non repellunt. Si sunt canones secundum spiritum: repel lunt. si sunt secundum carnem: nichil dico. nam parum faciunt ad Christum. Nam caro aduersus spiritum pugnat. Et vtinam carni numq̃ faueretur contra spiritum. spiritus enim verus sluminosus vnus syncaerus purus atque sanctus: quantuncumque viae carnis sint ĩpurae, viae Pauli: viaee spiritus. et canones id est regulae apostolorum: regulae spiritus. Ergo et bigamos et adulteros et fornicatores: censebimus episcopalis gradus celsitudine haud satis idoneos, integrae igitur carnis sint (et non diuisae: quales sunt et virgines et monogami. Verum quis negligentem et torpentem in specula collocabit enim quis ciuitati pateat incursus) exploratorem? quis desidi et somniculoso pastori: tute committet suum gregem? Plusquam ciuitas con mittitur explorationi et inuigilantiae episcopi et plusquam grex pastorali curae eius. Ergo qui vigil non esti attentus et admodum diligens: huic officio(huic bonoi et sancto operis non est committendus. Sit ergo vigil: sit et castus. Nam et pastores qui non sunt casti: minus inuigilant gregi. et qui positus est ad speculandum: si castus non est ciuitas in periculo versatur. Ergo episcopus: tum ne violet diuinae vnitatis simulacrum in magno illo Christi et ecclesiae sacramento tum ob vigilantiam castus sit. Sit ornatus, ornamentum animae: virtutes sunt et gratiae gemmis omnibus preciosiores. virtutes: quodam modo nobis comparamus. gratiae autem: nobis ex diuina luce infunduntur. et vt gemmae in tenebris: sic nostrae virtutes sine superna gratia. vt gemmae in radio solis: sic virtutes diuinarum grattarum illustratae radiis. Ergo episcopus huiusmodi fulgenti mentis ornatu debet esse conspicuus: debet esse ornatus. Sed et exterior ornatus deesse non debet: qui interioris preferat habitum. Attonsus crinis omnem superfluitatis et carnalitatis refectionem prae se fert et mundities exterior munditiem interiorem. Suparus ex candidissimo lino: castitatemt corporisque  ac animae munditiam ac sanctitatem. Sed ornatus exterioris singula percutrere: nimis operosum et huic forte loco minus accommodum. Sit igitur episcopus: mente ornatus: neque exterior ornatus ornatum mentiatur interiorem. alioqui est paries dealbatus: et sepul chrum adornatum. Nos omnes dum viuimus peregrini sumus: optantes tandem ad coelestem peruenire ciuitatem. iccirco nostrae conditionis memores: debemus esse hospitales et peregrinos libenter suscipere propter Christum qui voluit pro nobis et in hoc quoque  mundo esse peregrinus (vt nos tandem in aeternitatis hospitio susciperet. Et maxime omntum decet episcopum esse talem: vt qui Christum proxime repraesentet. Debet esse docilis. Nam si ebes esset ingenio et non facile legis intelligentiam capesses:non esset idoneus qui in loco illo sederet ad quem spectat vniuersae legis cognitioi non carnali quidem sensu et crasso sed spirituali et puriori examinanda. Sit ergo docilis siue docibilis id est qui facile doceatur: ettam desuper illucescente diuinitate proprio labore sacrarum rerum intelli gentiam penetrans. Non sit vinosus et in conuiuiis debacchans. nam irrepraehensibil enim vigilemi castumt ornatum hospitalem docilem esse oportet. Haec autem vinositas et (vt sic dicam) conpotationum frequentia et debacchatio: cunctis his aduersaturi et claudit diuinarum et humanarum rerum intelligentiam. Sit igitur sobrius: et quoad fieri potest abstemius Non sit percussor. At quis passorem gregis percussorem et mutilatorem deligeret? tanto minus debet episcopus humani gregis pastor percussor esse pugnator et mutilator. est enim vicarius pastoris aeterni: qui pascit non mterimit. qui contrita sanat: non vulnerat, qui author pacis est: non belli. qui luxata redintegrat: et omnia in se vnit dulci fouens ouem in gre mio sacris etiam suis humeris si deerrat reducens ad gregis vnitatem, non mutilat: non dis sipat. Talis ergo sit episcopus pastor: sub pastore illo constitutus lenis mansuetus pacisi cus, non solum percutiens factoi sed ne quidem verbo: quae dura est et saucianssnon sanans linguae percussio. pro quo et infra Paulus ωςρεοφὐτέρω μὴ ἐττιπληίfησ. quod interpretari solent seniorem ne increpaueris. Hic autem dicit μù τλήκτuη, quare et hic locus: ad percussionem etiam quae lingua fit tranfferri potest. Non sit turpis lucri sectatot. non decet episcopum turpis lucri esse sectatorem: neque a grege sibi creditoineque ; ab aliis more institorum et danistarum. Nam aeternus pastor dat: et omnia gratis datlet nonaccipit. Si quis autem episcopus pecuniam suam trapezitis tradit vt nummus nummum pariat si suis ad foenus traditisi mercaturas exercet: nomne turpe lucrum sectatur? Q si tur pe:etiam iniustum. iustitia eni: pulcherrimum animae ornamentum, quomon igitur hoc orna\pend
\section*{X }
\section*{1 Timo. }
\marginpar{[ p.205 ]}\pstart mento carens et turpitudinem loco eius induens: ornatus erit vt decet episcopum? Non igi tur sectet turpe lucrum: sed contral liberalis sit. Qui et cognitus sit propriae familiae bene praefuissel filios (si quos habuerit in subiectione et omni probitate ) ĩstituisse. Nam qui propriae familiae praeesse nescit: multo magis nesciet ecclesiae praeesse. qui enim nescit regere paucos: quomodo reget multos? qui nescit domum: quomodo ciuitatem? Ad tempora vsque Gre gorii septimi qui ordinis fuit Cluniacensium: adhuc sacerdotibus et diaconis licebat virginem habere vxorem. quae si prior moriebatur: non licebat secundam superducere. si autem superuiuebat: vidua manens sacros non contaminabat hymenaeos. Apostolicum nuptiarum ritum retinuerunt Graeci.:neque mutare voluerunt. agamiam acceptauerunt aliae ecclesiae. vnde plurimi per deteriorem icontinentiam lapsi: in pedicas inciderunt diaboli. Araneos videsi veneno turgentes: tam subtilia retia texere quae oculos fallere possint. quicquid incidit mortifero morsu necant: et primum quod aggrediuntur caput esti sensus auferentes. hae daemonis et laqueorum eius mira arte et subtilitate nobis praeten sorum et venenatorum morsuum: quaedam sunt adumbrationes. qui etiam retia plerumque  nectit: in iis quae clariora et sanctiora videntur. Caeterum qui praeficitur episcopus: non sit neophytus ad honores. nam qui nunq̃ administrationem gessit: nisi pedetentim probatus ascendat plerumque superari solet a gloria. cum nouicius est et venenifer araneus rete texit et flat superbiae ventus: incidit in pedicas eius. licet id in paruis rebus depraehendere. Nam qui doctorali laurea donantur: cum nouicii sunt svix seipsos depraehendant. sic eos quidam (non omnes dico sed quosdam) inanis tumor occupat. superuenienti bus autem annis: tumor ille paulatim recidit, vt cum iam sunt veteranil et plurimus eis honor iure debeatur: tunc nichili se aestiment humilitate ornati quos in nouiciatu inanis deformauerat elatio. Ergo non deligatur ad eam sublimem intendentiam nouicius. Neque satis est in enim loco testimonium bonum: si in aliis malum habet. nam etsi ipse forte in rete non incideret diaboli: hac tamen occasione vel exteri vel suorum aliqui inci dere possent. Araneus enim illic serpens ille venenatus antiquus: fila sua tendit. Ergo quoad fieri potest is deligatur huic prouiciae et sacrosanctae functioni committatur: qui ab omnibus bonum habet testimonium et omnibus gratus est atque acceptus. vt enim lumem solis omnium oculis gratum est: aeque et vera virtus omnium mentibus etiam plurimorum eorum qui aegrotant. Vt doctor gentium non ab hominibus hac laureal sed a Chri sto donatus docuit episcopos: subinde docet et diaconos qui ministri sunt episcoporum  et ex qbus si pro dignitate in ministerii sui functione se gerunt legitimo calle assumuntur episcopi. Quapropter summa authoritate dempta et prima sacrorum facultate: eadem ferme et horum et illorum officiat suntque vtrisque plurima communia. sed hi i imo et ministrando: illi in summo et praesidendo illa habentes. Et ideo subiungit Paulus de diaconis et familia eorumleorumque vxoribus hanc doctrinam: dicens. ¶ Diaconos similiter castos es se sermone non duplices (non multo vino deditos non turpis lucri sectatores habentes mysterium fidei in conscientia pura. et ipsi quidem probentur primum: deinde ministrent  cum sine crimine sint. vxores simili modo castas esse oportet. non calumniatrices vigiles  fideles in omnibus. Diaconi sint vnius vxoris viri: qui liberis honeste praesint et propriis domibus. Nam qui bene ministrant: gradum bonum sibi vendicant et multam fi duciam in fide quae est in CHRISTO IHESV. ¶ Ad institutionem episcopi septen decim annumerauerat. irrepraehensibilitatem monogamiam vigilantiam castitatem ornatum hospitalitatem docilitatem non vinositatem non percutiendi procliuitatem non turpis lucri sectationem mititatem pacabilitateml liberalitatemi bonam in re familiari praesidentiamlliberorum in omni honestate subiectionemi non nouiciatum testimonium bonum et ab iis quibus praeficiendus est et ab exteris. Ad diaconorum vero institutionem nouem enumerat. Primo. quod casti sint. hoc: carnis prohibet diuisionem. et tanto magis spiritus: quam facit idololatria et alterius q̃ veri dei cultus. Secundo, quod non sint mendaces. quid enim aliud sunt sermone duplices: nisi mendaces? sint ergo diaconi veraces: vt veritatis ministri esse possint qui Christo dominus est. Tertio. non sint vinosi. eos enim qui multo vino dediti et se amplius q̃ par est vino ingurgitant et proluunt: vinolentos et vinosos appellamus ac debacchatores. quod cum semper turpe sit: maxime tamen turpe cum cauponas non abhorrent vbi praebetur publicum vinolentiae et debacchationis exemplum. Igitur priuatus siue publicus sit bacchator id est multo vino deditus: non est ex praecepto apostoli immo Christi a quo doctrinam hanc cepit minister sacris altaribus assumendus. Sit ergo sobrius et abstinens aut saltem cibo potuque  temperans: quicunque huic officio est designandus. Quarto. Non sit sectator turpis lucri. hoc: prohibet omne lucrum ihonestum. vt quodcumque esset ex ludol quodcumque ex foeneratios  \pend
\section*{COM. }
\marginpar{[ p.3 ]}\pstart ne ex mercatura. Et si mercatura aliquiby possit esse honesta: non tamen diaconis. hone stum tamen lucrum: non prohibetur. vt quod est ex propriis fundis. Nam id bonoruni incrementumt a der benedictione et prouidentia venit:qui liberalis est in omnes-est itaque  hoc lucrum liberale: quod accipi debet cum multa gratiarum actionel ad liberalitatis opo ad rem ecclesiasticam si oportet ad famili am ad pauperes ad hospitalitatem id est ad aduenas prae se pietatem ferentes ad Christi similitudinem suscipiendos. non quod diaconi sit terras colere et arbores plantare. haec enim cura fidelibus idiotis committenda. Q uid si alicui honesto operi manuali intendunt et inde lucrum ad victumi vestitum et liberalitatis opus sectentur (Neque id prohibetur. Etenim Paulus maior omnibus episcopis maior diaconis omnibus: opere manuali operabatur id nobis relinquens vitae deo beneplacentis exem plum. Non igitur diaconus turpe sectetur lucrum:neque etiam tam sectetur honestum et liberalel q̃ vt eo habito liberaliter vtatur, nam qui haeret coelo coelestibusque ministeriis: non amplius terrae debet haerere sed prouidere. Nam et angeli prouident terris: sed non haerent terris sed patri coelesti cuius faciem semper vident. Quitum. Sint conscientia purai sacramenta dei ministrantes. Pontificis est: sacra mysteria celebrare, diaconorum est: ministrare et rebus pactis illa dispensare. at nisi habeant conscientiam puram: quomodo sacra mysteria rite dispensabunt? hoc igitur: ornatum iubet maxime interiorem qui in purita te conscientiae consistit. Sit igitur diaconus omnis: puritate conscientiae adornatus. Sexto. Non sit nouicius. Nam qui longo vsu aliqua in re probati non sunt: nouicii sunt. probetur igitur eorum idoneitas priusq̃ diaconatus donentur honore. non enim est honor paruus mensae domini ministrares pascere prius euangelica lectione fanctum Chri sti populum et deinde ipse pane vitae aeternae pastus: eundem vere coelestem panent immo supercoelestem et supersubstantialem propriis manibus reuerenter populo porrigerer quod est diaconi officium. Quod igitur iubentur prius probari et bene probatos assumi: id et nouicios excludit et irrepraehensibilitatem praecipit. tunc enim cum neq nouicii sunti neque aliquo vicio repraehensibiles: in sancto ministerio poni possunt. Septimo. Diacons si virginitatem non praeelegit: vxorem castam habeat non conuiciatricem sed vigilent et om nibus fidelem qualis decet sanctum ministrum. si carnem diuisit: casta non censeturl etiam si cum altero viro in nuptiis. tanto minus: si id accidisset per fornicationem, ergo quam licebat ei qui diaconus futurus erat accepisse vxorem: virginem esse oportebatiet sanctis moribus potius q̃ ampla dote adornatam. auaritia: infoelix et ifausta pronuba. castitas autemt omnibenigniloquentia sedulitas in re omni fidelitas: mulieris amplissima dos est. sine qua qui futurus erat diaconus et tanto minus qui futurus episcopus: nunq̃ virginem admittere debuit ad thorum. fortuna autem vltimum locum et abiectissimum semper teneat:virtus supremum. vt coelum summum locum tenet: et opes ex abditis effodiuntur terrae cauernis vt ex vltimis rerum locis. hic fortunae sedes: hic eius regia. Esto igitur diaconi vxor tali probitate isignita. Octauo. Diaconi si non virgines: solum sint monogami. nam digamum non licet assumi in diaconum: tanto minus trigamum aut polygamum. neque Christus nisi vnam ecclesiam: neque ecclesial nisi Christum vnq̃ habet. vnus sponsus: et sponsa vna. et tales decet in imitatione esse ecclesiae ministros. Et qui praesunt: non credant se esse sponsos. non sunt sponsi: sed vicarii sponsi. immo neq Petrus sponsus i neque Paulus: sed amici sponsit et fideles aeconomi eius. Et qui presunt et be ne praesunt: cum Paulo nomen sortiuntur aeconomi. sponsus vnus est: propter quem om nes sponsi et sponsae a pricipio mundi in imagine currunt. Et beati qui imaginem coelestis epithalamii in laudibus angelorum: non violant. Nono. enim liberos si habeant et rem familiarem non norit probe dispensare: idigni sunt qui ad ministeria dei assumantur dispensanda. Si ad tempora nostra respicimus: magna est in ecclesia dei indignitas. et timendum certe tantumdemque verendum: ne multi qui hoc nomine censenturi non sint dei ecclesial sed habeant libellum repudii. est tamen semper et erit sponsa Christi: sanca et immaculata. Satagamus igitur vias sequentes spiritus: illi iseri. et quicquid lesum fractum aut contritum fuerit: sanabitur reditegrabituri consolidabitur. Doctrinam et in iis qui praesunt et in iis qui subsunt imitemur Paulil et non simus filii discholiae: et omnia in pulchritudinet et mirabili quidem pulchritudine erunt. Quod annuat deus superbonus et supersummus etiam nostris temporibus: et si id non ad perfectumi saltem praeparationes iactamentaque videri. Cur de episcopis et diaconis scripserit Paulus ad Timotheum: causam aperit. quod forte adire non posset Ephesium ad episcopos et diaconos per loca istituendos: qui eo forsitan absenteristituendi erant. vt si ipse Timotheus aliquos sufficiat: studeat tales preficerelaut i mi ,nisterio ordinare. Ait itaque . ¶Haec tibi scribo: spe fretus quod ad te ocyus veniam. Si autem  tardauero: vt scias quo pacto oportet in domo dei conuersari quae est ecclesia dei viuëtis  \pend
\section*{X }
\section*{1 Timo. }
\marginpar{[ p.206 ]}\pstart columna et sedes veritatis. Et pro confessol magnum est hoc pietatis mysterium: deus ma nifestatus est in carne iustificatus est in spiritul visus est ab angelis praedicatus est in gem *tibus creditus est in mundol assumptus est in gloria. ¶Vt Moses insigniorem gratiam ha buit q̃ septuaginta duo qui de eius acceperunt spiritu: ita Paulus insigniorem q̃ Timotheus et alii coadiutores eius. Nam electus est: vt vniuersalis architectus. alii autem: vt cooperatores. Et Paulus propheta et Timotheus propheta:sed insignior Paulus ad co gnoscendum per spiritum sanctum qui idonei essent vt praeficerentur episcopisaut eligerentur ministri. non enim experientia exteriore egebant:qui cognoscebant interiora et penitissima cordium non ex sel sed ex spiritu sancto. Cum igitur per absentiam instruit Ti motheum: gratia maior instruit minorem et maius lumen allucet minori. Episcopus non praeficitur vt praesit terrenae domuilneque diaconus eligitur vt seruiat in terrena domo: sed ille vt praesit vt oeconomus in aliqua parte ecclesiae dei quae est spiritualis domus et domus dei viuentis et hic vt seruiat in eadem domo. Et nunc ecclesiam columnae comparat: sed columnae spirituali a terra ad summum vsque coelorum protensaelin cuius culmine sedet veritas qui est Christus super omnia benedictus. Haec est ecclesiae: in similitudine repraesentatio. Sed cum dixit de ecclesia quae sponsa est: subiungit de sponso qui sedet in columna. nam ille deus est: qui in carne manifestatus est. verbum enim caro factum est. Iustificatus in spiritu. requieuit enim super eum (etiam inseparabiliter) spiritus dominil spiritus sapientiae et intellectus Ispiritus consilii et fortitudinis spiritus scientiae et pietatis? et spiritus timoris domini. et ex plenitudine iustificationis eius: accepimus omnes. nulIns enim potest iustificari: nisi a iustificatione eius. vt nulla diurna lux:nisi a fonte folaris lucis. Visus est angelis. Nam ab angelis annunciatus. et natus in carne: q̃primum ab angelis euangelizatus. Inquit enim pastoribus claritate dei circumfulsis angelus praeemi nentior multa turba angelorum stipatus: nolite timere. Ecce enim euangelizo vobis gaudium magnum: quod erit omni populo. quia natus est vobis hodie: faluator qui est Chri stus dominus in ciuitate Dauid. et vbi dedit eis signum infantis ante omnia saecula et tunc in tempore nati qui fecerat et tempora: subito facta est eum angelo multitudo militiae coelestis laudantium deum et dicentium gloria in altissimis deo et in terra pax  in hominibus bona voluntas. Praedicatus in gentibus, nam et ipse iusserat apostolis: ite (inquit) et praedicate euangelium omni creaturae. et in omnem terram exiuit sonus eorum: et in fines orbis terrae verba eorum. Creditus est in mundo. sicut in prophetia scri ptum est. Populus quem non cognoui seruiuit michi: in auditu auris obediuit michi. Assum ptus est in gloria. Et videntibus (inquit) illis eleuatus est: et nubes suscepit eum ab oculis eorum. Et Stephanus primus testis et corona martyrum: vidit gloriam eius. Ecce (inquit) video coelos apertos: et filium hominis stantem a dexteris virtutis dei. Hic est qui sedet super columnam firmissimam et inconcussibilem vsque ad diem iudicii: cum mouebuntur virtutes coelorum. Et haec vniuersa magnum et ieffabile sunt circa nos pietatis eius sacramentum. quem introeuntem in mundum adorauerunt angeli dei: et quem nunc et semper adorat omnis exercitus coelorum. Cui sit honor et gloria super omi m eminens et glo riam et honorem: in omnia saecula saeculorum. Amen.  \pend\pstart ¶EXAMINATIO nomnullorum circa literam. ¶ Vulgata aeditio. ¶ Si quis episcopatum desyderat: bonum opus desyderat. ¶ Malui hic dicere episcopen et vti verbo Pauli: vt melius hic superintendential et officium episcopi intelli gaturi non idipsum cui commit titur. Stulti enim homines solent nunc quaerere, quanti hic vel ille episcopatus: attenden tes temporales fructus qui a subditis percipiuntur. Qz si intelligunt episcopent et eius ꝗ bonus est episcopus spiritualem fuperintendentiam tanti esse: similes sunt Simoni magot aut infoelici Ananiaesexistimantes spiritus sancti donum pecunia aestimari possei sed valor ille inaestimabilis est. Si vero negligens est episcopus: episcopatus nichil valet. quomodo enim superintendentiam aestimares:quae nulla est? quod certe nichil est: aestimari non potest. Terrena igituri quae bene intendentes sequuntur: non valori non precium est illius sanctae superintendentiae sed quaedam spiritualium operum hypostasis vt sensibilis mundus hypostasis spiritualium et ĩtellectualium rerum et vmbra lucis hypostasis. Qui bene praesunt vigilant superintendunt: lucem ministrant et in modica subditorum vmbra capiunt refrigerium. Secus est de iis qui suae vicis functionem non intelligunt. volunt enim hi in sola temporalium vmbra requiescere: nichil curantes summi patrisfamilias gregem per suam secordiam mille modis interim (dum torpent in vmbra) interire. hi episcopent hi episcopatum numq̃ desyderauerunt: sed solam episcopes solamque ipsius episcopatus vmbram quam falso nomine episcopatum appellant. Paulus. ειτισ ἐττις σuowη̃σ ὀρέγεται, καλου ἔρνου ἐπιθυμει. Est enim e̓ πισηοιη intendential  \pend
\section*{COM. }
\marginpar{[ p.3 ]}\pstart ,opus officiumque episcopi. ¶ Vulgata aeditio. Vnius vxoris virum sobrium ornatum pru» dentemi pudicumi hospitalem doctoremi non vinolentumi non percussoremi sed modestumi non litigiosum non cupidum. CQuod hic vetus interpres dicit sobrium: graece habetur puφά Aiop: quod magis vigilem significat. Et quod dicit prudentem: graece ha betur σῶ φρoμαι quod alibi sobrium dicit. nos autem modestum et castum et temperatum: potius dicimus. et σωφρωp ad significandum qρópιμopid est prudentem: id non nisi metaphorice fit quod temperantia prudentiae soleat esse comes prudentisque conseruet saluetque iudicium. Et debuit hic secundum locum tenerel ponique ante ornatum. Et hic ornatus mundiciem dicit in cultu exteriori. Non est enim episcopus sponsus: sed vicarius sponsi sponsumque repraesentans. Decet igitur ipsum in habitu mundo castoque oculis castae et virginis sponsae se semper praesentare. Pudicum: hic abundat. Quod vetus interpres dicit doctorem: potius designat docilemi abilemque ad docendum. quomodo enim alios docere posset: qui ne quidem doceri posset ? Vinolentum siue vinosum eum intelligimus: qui deditus vino propemodum bacchatur. dictum apud Graecos: quod senper iuxta vinum assideat. Quod cum omni Christiano: maxime episcopo turpe esti ac vitandum. Post percussoremi deest vocabulum vnum: quod Paulus dicit ἀι σ χροxepAh̃. nos si vsus admitteret: enim composito vocabulo diceremus turpilucrem.verum quia apud latincs parum vsitatum est: circumlocutione magis vtendum est et dicendum turpis lucri cupidum aut turpis lucri sectatorem. Et quod vetus interpres modestum di cit: mitemi et clementem significat. Quod non cupidum: cupidum specialiter designat auarum argentique cupidum. Paulus. puφά λιορ, σω φρορα. κόσμιορ, φιλόfεpopAιΑαΚτιΚόμς μη τταροιρομ, μἠιττλήRτημ, μὴι ἀισxpοkςρAῆ, αλλέSιEιRù, ἀμαχορ, ἀφιλάρτυροp. ¶ Vulgata aeditio. Diacones similiter pudicos non bilingues. ¶Dicendum: diaconos. nam diacon non dicitur: sed diaconus. Paulus. λιακόνουσ ωσ»αύτως σεμυους αη λἰλόγουσ. ¶Vulgata aeditio. Sobrias fideles in omnibus¶Dici et hic potest: vigiles vbi dicit sobrias. Paulus. puφaλέουο, πίςας ἐy τασι. * ¶ Vulgata aeditio. Diacones sint vnius vxoris viri: qui filiis suis bene praesint. ¶Dicen dum diaconi. Et suis: abundat. Et cum diciti diaconi sint vnius vxoris viri: deliro esset et vanissimus sophista importunusque garrulus si id coniunctim intelligeret et non seorsum ac diuisim. ac si diceremus: diaconus quisque sit vnius vxoris vir. Tales enim obliquo mentis oculo omnia cauillantes a spiritu tam procul absunt: q̃ sathanas a sanctorum coetu angelorum lusca monstrat aut penitus exoculati Polyphemi. Paulus. AiάRopoι ἐς σοcαp μιᾶσ ΤυμαιΚὸσ ἀμAρεσ, τέπρωμ Καλῶς προίξάμεpοι. ¶ Vulgata aeditio. . Haec tibi scribo fili Timothee:sperans me ad te venire cito. ( Particula fili Timothee: abundat. Paul. ταῦτάσοι γράφω, ἐλωιτωμ ἐλοεῖμν ττρὸς σὲ τάχιον. ¶ Vulgata , aeditio. Quae est ecclesia dei viuu columna et firmamentum veritatis. ¶ Potius sedes aut fundamentum q̃ firmamentum dicitur. Verum et sedem dicere malui q̃ fundamentum. ne putetur hic dici fundamentum:vt et cum dicitur. Fundamentum enim aliud nullus potest poneres praettr id quod positum est: qui em Ihesus Christus. nam hoc Christo conuenit non ecclesiae: quae est fundata et superaedificatalnon fundamentum. Christus autem ecclesiae fundamentum. spiritualium enim fundamenta sursum sunt non terrenorum mo re deorsum (fundamenta inquit eius in montibus sanctis) et fundamentum quod maxi me per eminentiam dicitur super omnes est. Cum igitur ecclesiam comparat columne: ea quae in terra est non est principiumlsed finis et extremum. quae enim fundamento sunt propiora: et illa sunt priora. Vt supremae virtutes coelestes proxime in fundamento posi tae sunt et montes sancti in visione spiritus: in quibus et sedet deus, magis igitur consentaneum visum est nec ab re quidem: vt sedem potius interpretaremur q̃ fundamentum. id enim magis ecclesiae conuenit: vtpote in cuius summitatel vt in summitate columnae a terra ad coelorum summa eleuatae sedet deus. Paulus. Hτισ εὐσὲμ ἐ λελησία 9ςοὺ φωμν  τoσ, 5ὐλοσ Και εὐλραίωμα τῆς ἀλμοςίας. ¶ Vulgata aeditio. Et manifeste ma* gnum est pietatis sacramentum: quod manifestatum est in carnel iustificatum est in spiritul  apparuit angelis praedicatum est gentibus creditum est in mundol assumptum est in glo» ria. ¶ Potius pro confessoiq̃ manifeste: dicendum. significatque id pro explorato et certot omni propulsa dubietatis caligine. Atque dicendum deus manifestatus: non quod manifestatum em. et caete ra sequentia adiectiua neutra l’in masculina conuertenda. dicendo: iustificatus est praedicatus esti creditus est (assumptus esti vt omnia in suum substantiuum deus referantur. Paulus. καὶ ὁμολογουμέρωσ μέγα ἐσὶ το τῆς ἐυσεφέιασ μυςήριοp, θεὸσ ἐφα- pερώοη ἐμ σαρΚι, ἐ λιΚαιώομ ἐρ ττμςύματι. dιφοη ἀτγέλοις, e̓kHpυχOu e̓μ topscip, exiseuou en κόσμω, ἀpελήφoη ἔν αόςα. Nunc ad sequetia transeüdum.  \pend
\section*{X }
\section*{1 Timo. }
\marginpar{[ p.207 ]}
\endnumbering\beginnumbering\section{CAP. IIII.}\pstart \huge\textbf{H}\normalsize Oc facto magnoque pietatis reserato sacramento: prophetat Paulus Timo theo de iis erroribus qui futuri erant circa ciborum abstinentias et nuptiarum quasi haec per se mala essent. Qui error postmodum apparuit: qui fuit Encratitarumi et eorum quos appellarunt Cataphrygas. et expe diebat Asiaticos prophetia praemuniri. quia error ille:apud eos plurimum inserpsit. Ait ergo Paulus. ( Spiritus autem eloquibiliter dicit: q in posterioribus temporibus apostatabunt nonnulsi a fide attendentes spiritibus erroneis et doctrinis daemoniorum eorum qui simulatione ,mendacia loquuntur cauterio inustam propriam habentes conscientiam prohi,bentes nubere iubentes abstinere a cibis quos deus creauit ad perceptionem cum ,gratiarum actione a fidelibus et iis qui cognouerunt veritatem quod omnis creatura .dei bona est et nichil excludendum quod cum gratiarum actione sumitur. sanctifica,tur enim per sermonem dei et locutionem, ¶ Haec est prophetia de Encratitis et de haeresi Phrygianorum. Quis enim nescit delirum Montanum aliena loquentem Priscillam et Ma ximillam insanas vates et alios aut authores aut fautores horum errorum: spiritibus at tendisse erroneis et innumera mendacia loquutos fuisse pietatem simulantes conscientiam habentes vlceratam et diaboli cauterio penetratami ne cicatrix aboleri posset? Qu si quis amplius cupit has tenebras (quas vult depellere apostolus) cognoscere: Epipha nium legat, oculus tamen qui est amator lucis: his minime cupit inuersari tenebris nisi cum necessitas vrget ad illas depellendas et ad verae lucis reparationis diem. alioqui timem da est curiositas: et ex curiositate quaedam lucis purius admittendae impraeparatio et ingressus maligni facilior admissio. Quid igitur mandat Timotheo? Vt praedicet: et om nem creaturam esse bonam quam deus condidit et omnem escam rationabilemlbonam homini essenullamque excipiendam quae cum gratiarum actione percipitur et nuptias esse bonas. venturosque esse: qui aliter ex doctrinis daemoniorum quae perniciosae haereses sunt doceant vt a vera creatoris pietate abducant, etenim cibi omnes rationabiles et nu ptiaes per domini sermonem sanctificantur. Rationabiles dico. nam serpentes virosa et humanae carnes et multa id genus licet creaturae bonae sint: non funt tamen rationabi lis cibus. neque nuptiae rationabiles: si mulieri iungatur eunuchus aut mulier viro qui puertere naturam vult propriarum nuptiarum vsum ludibrio habens. nam neque illi cibil neque hae nuptiae: sed feritates et circa naturae cibos et nuptias pessima daemonum ludifi camenta quae sanctificationem non recipiunt. omnia enim sanctificationi apta: ordinata sunt. Subdit ergo Paulus ad Timotheum. ¶Haec admonens fratres: bonus eris mi,nister IHESVCHRISTIt enutriens sermonibus fideil et bona doctrina quam assecutus ,es. (LDat insuper ei alias doctrinas quae cuique volenti alios istituere et in pietatis officiis seipsum exerceret necessariae sunt. Prima vt cum aliorum suscipit institutionem: propha nas fabulas et aniles non immisceat. Praebent enim eloquia et agiographorum castorumque scriptorum dicta (si quid sermoni inserendum est) amplissimum campum. Ait itaque - ¶Prophanas autem et aniles fabulas deuita. ¶ Certe in institutione aliorum: prophana inquinantiai et anilia omnino verba deuitanda sunt. et non solum publice: sed et omnino priuate vitanda sunt. verbo domini: ornatus mundi prodiit pulchritudoque syderum. fic quicquid ab ore nostro prodit qui sumus filii verbi dei: ornatum pulchrumque  esse debeti et iternae virtutis et gratiae lucel vt micantibus stellis respersum. Sed nume cum neqssimis fabulis cum mollibus elegiis cum lenonibo lenis meretriculis parasitis comoediarum ita ab ineunte aetate emolliuntur adolescentum animi:vt cum ad adultam aetatem veniant  nichil placeat nisi prophanumi nisi molle et intemperans nisi fabulosum et quod deliriis vetularum magis indignum sit ac perniciosum. Christus suos authores habet: suos hymnidicos et epicos et (vt sic dicam) σεμpοτroιhτασ id est semnopoetas omnes suos historicos suos philosophos et sapientes. in quibus qui instituuntur: castas mentes repor tanti luci spiritus sanctif et sermonibus eius peruias donec aeternus eis oriatur dies non ante pubem inille modis corruptas et incestas per molles fabellas veneres cupidines  comoedias satyras in plerisque ita turpia turpiter repraehendentes ac si quis ante virgines castitatem laudans (quod etiam nunc dicere pudet) denudet pudendas ac si veneficia culpans ante quos culpat venena misceat. et sit veneficus: qui ante veneficii culpam quid esset veneficium prorsus nesciebat. hae fascinationes daemonum sunt haec prophana: et voce Pauli prae conis altissimi deuitanda et procul ab omni gymnasio Christi arcenda, illa enim sunt de gymnasio idolorum : quod ei aduersatur quod Christi est. Se, cunda quam hic ponit doctrina: de exercitio ad pietatem esti et est haec. ¶ Exerce au\pend
\section*{COM. }
\marginpar{[ p.4 ]}\pstart tem teipsum ad pietatem. Nam corporalis exercitatio: ad paucum vtilis est. pietas au» tem ad omnia prodest: promissiones habens vitae cum praesentis tum futurae. ¶ Exercitatio corporis motio corporis quae sanitatem viresque cum accommodata et temperata fuerit conseruat. Sed pietas: est mentis ad deum conuersae motiolquae excessum non habeti vitam conseruans et praesentem et aeternam. haec:ad aeternitatem valet. illa: solum ad praesentem vitam. praesens vita si ad futuram conferas: punctum est ad lineam infinitam. corporalis igitur exercitatio, quae solum ad praesentem vitam pertinet: ad paucum vtilis em, sed exercitatio mentis secundum pietatis motionem: ad multa imo ad omnia vtilis est. Nam et ad praesentia et ad futura quae omne bonum continent. Et piis: per os sui apostoli promittit deus vtramque vitam et praesentem et futuram. mens pia et religiosa: oculus est semper ad solem conuersus omnia quaecunque facit ex illa luce hauriens iaut ad illam lu cem reducens. suntque omnia opera pii: lucidam et luminosa, quomodo enim qui ad solem conuertitur: quicquam facit quod non illuminetur Tertia doctrina vt speremus a deo salutem: et a re alia nulla. Nam ipse est saluator omnium hominum. quis enim in hoc mundo om nes homines saluat: nisi qui vitam omnibus hominibus largitur et prouidet vnde habi ta conseruetur? Sed maxime saluator est hominum fidelium. nam eos et in hoc mundo saluat per prouidentiam: et in altero per gloriam et vitam sempiternam. infideles autem solum in hoc mundo saluat: i altero vero (vt par est) gloria vitaque priuat aeternal per iustitiam. ergo prouidentia: communis est tam bonis q̃ malis. sed gloria et vita aeterna: solum propria est bonis. ignominia extremaque confusio et mors aeterna: malis. gloria inq̃ propria est bonis sed qui fideles sunt. aeterna autem confusio malis: sed qui infideles sunt. aut simpliciter: vt qui neque verbo quidem credunt. aut opere: vt qui verba fidei habent sed factis fidem impugnant. vt qui dicunt non fornicandum et fornicantur non furandum et furantur: horum cum infidelibus portio est. et q̃diu sic viuunt: nomine infidelium continem tur. si autem resipiscunt ad deum honorandum quem inhonorarunti ad amandum quem offem 19 derunt ad colendam iustitiam quam transgressi sunt (vniuersa enim dei mandata iustitia con tinet) nomen mutabunt eruntque rursus fideles ex infidelibus. Ait igitur Paulus. ¶ Fidelis sermo et omni acceptatione dignus. in hoc enim et laboramus et i properia sustinemus: » quia sperauimus in deo viuente qui est saluator omnium hominumi maximel fidelium, ¶Qui spem vniuersam suae salutis in deo collocat: laborem omnem et improperia mundi facile tolerat. quia sperat in eo: qui saluatlaut damnat. Agricola cum seminat siue serena dies sit siue turbulenta: neque in opere suoi neque in praesentis temporis aut serenitate aut aduersitate sperat (nam spes futuri est) sed in sola dei prouidentiai quae sola potest aut messem multiplicare aut eam si vult perdere. fecit tamen quod potuit: et vt agricola fidelis. Oz si agrum non arasset si semen non iecisset: non speraret. Simili modo qui bene operatur in rebus secundis: semen salutis iacit in serenitate dei. qui autem in aduersis et cum dura premunt: semen salutis iacit in turbulentia dei. sed neque in secundist neque in aduersis neque in operibus suis confidere debet: sed in eo solo qui facere potest vt iacta semina adolescant in sa lutem. non enim deus aliter dat salutem operibus: q̃ vitam mortuis in terram iactis seminibus. Quarta doctrina. Vt praedictarum doctrinarum luce caeteros illuminemus. et sic exemus A plum vitae in sermone conuersatione dilectione spiritu castitatel et fide nosipsos constituamus: vt nemo nos contemnere possit. et hoc etiam ab adolescentia discamus. Nam si ita vi2O uamus vt adolescentes nos nemo contemnat: tanto minus contempturi sunt cum fuerimns senes. Ait ergo apostolus. ¶ Praecipe haec: et doce. Facito nullus tuam iuuentam contemnat:  sed exemplum esto fidellum in sermonel in conuersationel in dilectionelin spiritusin fidel , in castitate. ¶ Quod praecipit Timotheo: omnibus praecipit maxime qui aliis vt praesint praefecti sunt. Nam praecipere et docere: superiorum sunt ad inferiores, neq enim ima illuminant summa: sed summa imis suae illuminationis radios imittunt efficiunturque ima per radiorum susceptionem (quodque pro captu) summiformia. sicque minores maiorum exempla sequentes: debent quoad possunt per imitationem maioriformes in sermone l in conuersatione i dilectione in operibus spiritus i fidel in castitate fieri. Haec enim a superioribus i ine feriores sic lucere debent: quemadmodum a sole omnia inferiora luminosa videntur suntque  luciformia. Quinta doctrina etsi Timotheo quoquo pacto peculiaris videatur: vniuer; salis etiam ex eam intelligi potest quae sit haec. Diuinae lectioni exhortationil doctrinae intem dere debemus donum dei quodcunque nobis datuin est (non negligere in illis immorari. Nam hoc facientes: nos ad dei salutem adaptamus et eos identidem quibuscum versamur et nos audiunt ad salutis portum conducimus. Non subet nos impurae intendere lectioni (procul absit Plautus Terentius Martialis Tibullus aut tenerorum lusor amorum et reuera melius illusor animorum) sed sancae. Sacrum pectus est: cui sanda in  \pend
\section*{X }
\section*{1 Timos }
\marginpar{[ p.208 ]}\pstart scripseris euangelia cui castam Pauli doctrinam, et quicquid doctrinae non aduersatur spiritus. si alium ex sancta lectione non sufficis exhortari: exhortare teipsum. si non sufficis alium docere: doce teipsum. donum dei non neglige nullus est: quem non dono aliquo mu neret. quod qui negligit: datorem spernere videtur. et hoc ipso: se donis aliis constituit indignum. Q ui his doctrinis immanet qui in his se exercet: ad pietatem se exercet quae ad omnia valet et ad propriam et ad aliorum salutem, et ad praesentem et ad futuram vitam. tres 21 itaque supremae doctrinae: secundam quoquo pacto aperiunt et dilucidioremsaptioremque reddunt. *Subiungit ergo Paulus ad Timotheum hoc pacto. ¶Donec veniam: intende lectionil » exhortationi doctrinae. noli negligere donum quod in te est: quod datum est tibi per prophetiam cum impositione manuum presbyterii. Haec meditaresi his esto: vt promotio tua » manifesta sit in omnibus. Attende tibi ipsi et doctrinae. immane illis. hoc enim faciens: et  teipsum saluabis et eos qui te audierit. ¶ Pres byteros qui et apud latinos presbyter: senior dicitur. presbyterium: functio et potestas seniorum erat qui manus capitibus initiandorum purgandorum aut perficiendorum imponebant sacra precamina fundentes diuinamque  implorantes opem. Qui cum manus imposuerunt Timotheo: donum accepit prophetiae donum intelligentiae prophetarum et sanctarum scripturarum. Et qui sanctas et impollutas manus ei impofuit: Paulus fuit vt ex sequenti libro cognoscetur qui et sanctos fudit precatus dona spiritus sancti a deo eidem obtinens. Verum nunc hoc nomine presbyter vtimur: non vt dicit aetatem sed ordinem. ita vt cui sit senium mentis et prudentia spiritus: et si quintum et vicesimum annum nondum sit consecutus modo sacra maiorum functione ad hoc sit consumatus pres byter seniorque dicitur. est enim is ordo: non initiatio sed consumatio quaedam. sed talis quae non debetur nisi iis solis: qui mente sunt canal cui maior debetur reuerentia q̃ caniciei secundum corpus. mysteria enim illa peragunt: quae omnibus aliis sunt augustiora. Et vt in terrenis qui summam et absolu tiorem tenent potestatem terreni reges sunt: sic et isti summam et absolutam diuinorum tenentes potestatem sacrorum reges. Viuit adhuc (nisi me fama fallit) sanctam et ĩpollutam (vt arbitror) vitam ducens vir quidam asperrimae vitaelcui puellulo per plures an nos angelorum ministeria circa diuinum presbyterii officium patebant. nunq̃ res sacra illi peragi visa est: quin coelestis quaedam corona supra presbyteri caput a consecrationis tempore tamdiu appareret dum domint corpus sanguisque sumeretur. tunc illa: inuisibilis disparebat corona. Hoc et verum credo: et huiusmodi virum adhuc Mantuam tenere. Sed quid hoc: nisi signum quod reuera regale est sacerdotium absoluta quaedam et mirifica potestas cui assistunt angeli dei et quam reuerentur? Ergo sacerdotes omnes: haec mediteminilin his estote. promotio vestra manifesta est angelis non solum hominibus. saluate vosipsos quoad in vobis est: et eos qui vos audiunt et qui vobis assistunt. in vobis autem est apparare aptare disponere: dei autem est consumare. Gratias ei agite peren nes: qui vos reges fecit ad opus immortale qui rex regum est et cuius solius est reges terrenos et spirituales creare. Cui sit honor maiestas decus et imperium: in omnia regna regnantium et haec in omnia saecula saeculorum. Amen.  \pend\pstart TEXAMINATIO nomnullorum circa literam. ¶ Vulgata editio. Spiritus autem mani ,feste dicit: quia in nouissimis temporibus. (¶Quod vetus interpres dicit manifeste: aduerbium esti loquibiliter significans ac si diceremus. spiritus autem loquens dicit: quia aduerbio parum vti consueuimus. Et non dicitur in Paulo i nouissimis per supellatiuum: sed comparatiuo vtitur. Vt etiam prophetia non multo post tempora Pauli magna sur gente haereticorum caterua a vera fide desciscentiumlimpleta esse potuerit. Paulus. T o λe ττρεῦμα ρuτῶσ λέτει, οτι ἐp υsέρoισ xαιpoi͂σ. ¶ Vulgata aeditio. Loquen tium mendacium. ¶ Ne loquentium videretur referri ad daemoniorum: satius visum est vt genitiuus ille et sequentes in rectos couerterentur vt aperte referat quod supra positum est quidam. tunc enim sermo latinus: omnino magis est legentibus peruius. Alioqui itelligatur pseudologon id est loquentium mendacium referri ad daemoniorum. loquuntur enim daemonia in simulatione sanctitatis et religionis mendacial quo simplices homines et humiles deiectique sensus magis decipiant. Nos primum sensum cepimus referendo ad hominum: vetus autem interpres forte secundum referendo ad daemoniorum. neque hunc neque illum improbauerim. Nam quae huiusmodi homines in simulata pietate loquuntur: ea ipsa sunt quae fallacia loquuntur daemonia. Paulus. ςυ λολόT Oμ.  ¶Vulgata aeditio. Et iis qui cognouerunt veritatem: quia omnis creatura dei bona est. ¶Potius dicendum quod omnis creatura dei bona. ne quia: confirmatiuum aut redditiuum causae videatur quod potius videtur explicatiuum. cognoscere enim quod omnis creatura dei bona est nichilque excipiendum quod cum gratiarum actione percipitur: est cogno\pend
\section*{COM. }
\marginpar{[ p.4. ]}\pstart scere veritatem. Paulus. Καὶ ἐπεγρωκόσι ἀλήθειαμ, οτι πὰμ Κτίσμα θεου. Is Vul ,gata aeditio. Sanctificatur enim per verbum dei et orationem. ¶ Potius: et locutioncmi aut supplicationem. Non enim hic ponitur vocabulum υτροσευχὴ quod passim oratio verti solet vt in principio secundi capitis hutus vbi dicitur. Obsecro igitur primum om nium vt fiant preces orationes. sed ponitur vocabulum: quod vetus interpres illic vertit postulationes. Paulus. ἀτιάδεται γἀρ λιἀ λόγου θεοῦ, καὶ ἐυτ εύθεεος. ¶Vul ,gata aeditio. Haec proponens fratribus: bonus eris minister CHRISTI IHES Vr enutritus verbis fideil et bonae doctrinae quam assecutus es. (¶Potius: haec admonens fratres. Et Ihesu Christi. Et enutriens dici potest:quod per praesens loquituri tametsi vocem pas siuam habeat. Et dicendum esset ex aliquibus codicibus graecis: et bona doctrina sexto casu. nam in graeco eiusdem casus est cum eo quod vetus interpres dixit verbis: non cum fidei, quare codem casu verti deberet. Paulus. ταῦτα ὐττοτιθέμεpος τοῖσ ἀAελquodomniσ, καλὸσ εὐση Αιάκομοσ ἰuσοῦ χριςοὺ, ἐμτρεφόμενοσ τοῖς λóγoις τuς ττίσευως, καὶ τῆ καλῆ Αιλασκαλία. Sunt etiam codices graeci:qui hanc vltimam particulam habent: Καὶ τῆσ καλη̃͂ς λιλασκαλίασ. Et tunc dicendum per genitiuum et bonae doctrinaei et sic vetus legit interpres. Quapropter et sic etiam legi potest. ¶Vulgs ta aeditio. Quia speramus in deum viuum. ¶ Potius dicendum: sperauimus per praeteritam. Et in deo viuoisiue viuente. Paulus. orι hλπίπαμεν ἐ ωι θςῷ φopTι. ¶Vulgata ,aeditio. Nemo adolescentiam tuam contemnat.(] Hoc prae ceptum non aliis sed ipli fit Timotheo. Ad quod magis inuendum opereprecium magis duximus sic esse dicendum facito nullus tuam iuuentam contemnat. Paulus. μηλείσ σοῦ τῆσ ρεότητοσ καταφρομεετω* ¶Vulgata aeditio. In conuersationet in charitate in fide in castitate. ¶Dicendum in charitate in spiritu. deest enim i spiritu post charitate. Paulus. ἐμ ἀμαςροφη̃, e̓μ αpαττη. νἐμ πτρεύματι, ἐμ ττίgει, ἐμ ἀγpεία. ¶ Vulgata aeditio. Noli negligere gratiam quae i ,te est. ¶In graeco vocabulum quod passim donum interpretari solet habetur non gratia et si in re et significatione ( vt dictum alias est)cum gratia conspiret. Paulus. 1iii ἀ- εμέλει τοῦ ἐν gοὶ χαρέσματοσ. ¶ Vulgata aeditio. Insta in illis. ¶ Illud insta designat immanere. Et respicere videtur id quod prius dictum est: hec meditaret in his estot immane illis id est mane in illis contine te in illis neque velis ad aliena excedere. Q uod praeceptum etsi in sacris eloquiis maxime est obseruandum: si tamen in caeteris disciplinis etiam humanis obseruaret meditando instando et immanendo iis quae a probatis authoribus tradita sunt paruo tempore omnes disciplinas acquirere possemus et ex illis legitimo quodam calle ad diuina transferri vt ex sensibili mundo ad ĩtellectualent ex humano spectaculo ad diuinum (vt sic dixerim) theatrum. Q quia probatorum dictasneque adole scentes neque adolescentum praeceptores hac nostra tempestate meditanturiin illis non sunt illis non immanent sed ad aliena abiguntur pascual magis infoecunda q̃ solitudines Librcaesq̃ Arabum auia: nunq̃ ne vnam quidem (vt par est) etiam humanam adipisei valent disciplinam. et cum se omnia scire praesumunt: egregie delirant. et cum culices venau sunt: arbitrantur se cepisse elephantes. Sed omnia apud eos oracula: leuiora sunt foliis Sybillae. Et vtinam saltem sacras (si dici volumus theologi) meditaremur scripturas il lis insisteremus illis immaneremus et quasi intra muros paradisi contineremur et non ad aliena raptaremuri quae prohibet apostolus. Quod donec fiat: cum in humanis tum in diuinis lux diuina (etsi nobis affulgeat) nunq̃ capietur. Valeant ergo omnia exoticas barbara foeda olida et quae quicquam impuri et foedi (quis elegantia sermonis adumbrati) continent: vt lux capiatur diuina. Tam enim ab elegantibus improbam etsi festiua oratione continentibus sententiam iq̃ a barbaris inconcinne et male perstrependo ad aliena trahentibus: cauendum est. Paulus. ἐττίμεμε ἐαυτοῖσ.  \pend
\endnumbering\beginnumbering\section{CAP. V.}\pstart \huge\textbf{I}\normalsize Nstituit insuper apostolus Timotheum: quomodo in omni officio attendat decorum( et personis se in omni honestate accommodet. Et ait. (Seniorem ne increpaueris: sed obsecra tanq̃ patrem: iuniores-tanq fratres. Anus: vt matres. Adolescentulas: vt sorores i omni castitate. viduas honora:quae vere viduae sunt. ¶ Si in eo qui praeest tantam in docendo et corrigendo requirit mansuetudinem et humanitatis graetiam: q̃ non intelligunt illi qui nullam habentes aut in personas aut in coetus superioritatis praeeminentiam tanta austeritate carpunt quasi ii ipsi sint qui mittant in gehennam putantes hoc ipso seipsos iustos esse sed discant formam docendi et mansuetudiuem ab apostolo: et quisque gradum suum cognoscere. Seniores obsecrent tanq̃ patres etiam si episcopi erunt vt Timotheus aequales tand̃ fratres lon ge etiam aetate inferiores aut vt fratres aut vt filios. Anus: vt matres. Adolescentulas:  \pend
\section*{x }
\section*{1 Timo. }
\marginpar{[ p.209 ]}\pstart vt sorores. videant tamen ne id faciendo: impuras sentiant carnis motiones, essent enin infirmi: magis forte et doctrina et medicina egentes q̃ quibus doctrinaml medicinamque pararent. id igitur in omni castitate et puritate fiat. quod nemo potest aut vix: nisi qui spi ritualis est. esse autem spiritualem: non cuiusq̃ est quanticumque se aestimetised est dei donum. Viduas quae vere viduae sunt: honorent. vere viduae sunt vt turtures carentes compare: quae nulli amplius mortali viro haerere volunt sed quod reliquum est vitae soli Christo haerere solique  placere simplicitatem in se gerentes columbinaml et castitatem seruantes turturinam. Verum quia de viduis locutus est: amplius prosequitur. nam quae viduae sunt: aut filios habent aut nepotes aut sunt solae relictael pariter orbatae viroffiliis nepotumque solatio. Nunc autem pro ea cui sunt filii aut nepotes: hoc ponit praeceptum. Filii aut nepotes domum cum pietate regant: et viduae matri mutuam prouidentiae vicem rependant. Sic ergo ait Paulus. ¶Si qua autem vidua filios, aut nepotes habet: discant primum proprias domos religiose curarel et mutuam vicem parentibus reddere. Nam hoc: coram deo acceptabile est. ¶Certum est hoc praeceptum tunc habere locum: cum filii aut nepotes tantilli sunt qui iam possint domum gubernare et exteriorem curam gererel vidua autem interiorem, vt quae dudum didicit intra domum esse materfamilias. et haec vidua: sed non vsque adeo vere vidua, nam loco viri: filiorum aut nepotum subleuamen ac solatium habet et plerumque circa rem familiarem perinde occupata viuit vt virum habens nisi nurus eam a cura rei internae leuent. Sed quae nam ,omnino et simpliciter vere vidua dici debeat: subiungit Paulus dicens. ¶ Vere autem vi,dua et quae sola manet: sperauit in deo ( et imanet orationibus et obsecrationibol nocte ,dieque . ¶Hanc ergo vult dici vere viduam: quae sola relicta sperauit in deol neque diel neque  nocte ab orationibus et obsecrationibus cessat. sperauit dico magis q̃ in viro dum viueret: et magis deo spiritu adhaeret q̃ adhaesit corpore viro. ideo orando magis deo et die et nocte famulatur q̃ vnq̃ terreno et mortali sponso. quae licet neque filiorum neque nepotum solatia quaerat: solatur tamen animo magis q̃ prolis aut consanguineorum vllo posset famu letio. quaecumque ; huiusmodi est: omnino simpliciter et vere vidua est. De ea autem que om , nino vere vidua non est: subdit apostoloi dicens. ¶ Quae autem in deliciis degit: viuens mor *tua est. ¶ Deliciae: carnem nutriunti et extingunt spiritum. vita autem carnis: mors spiritus est. mors autem spiritus: vera mors est. quare et huiusmodi vidua cum spiritu mortua sit: etsi carne viuens reuera mortua est. Sed ex doctrina viduarum mulierum: quid de viduis viris et veres et non vere viduis sentiendum sit edocemur. Coniugatorum vita: actiuorum est. vere autem viduorum et virorum et mulierum: contemplatiua esse debet. etsi non in summo saltem in imo: vbi oratio et depraecatio in humilitate et simplicitate collocatur. Iubet apostolus vt quisque suarum viduarum curam geratisiue pater filiam siue filius matremi siue nepos auam suami viduam habet. Ac ait. ¶ Et haec praecipe: vt irrepraehensi'biles sint. si quis autem suis et maxime domesticis non prouidet: fidem negauit et est *infideli deterior. ¶ Quare fidem negauit? Quia religio munda et immaculata apud deum et patrem haec est: visitare pupillos et viduas in tribulatione earuml etiam ersi visitantium nichil sunt. tanto magis igitur cura gerenda est suarum et domesticarum viduarum. qui igitur id non facit: fidem ac religionem negauit. Certum est facto fidem negasse. sed et si dicit ad se nichil attinere prouidere eis: fidem cum verbot tum facto negauit. et est infideli deterior. nam multi infidellum: suis prouident. et si non prouident: fidem non abiiciunti quam nunq̃ susceperunt. Quartum ponit praeceptum de viduis quibus ex ecclesia prouidendum et quae de mensa ecclesiae sunt pascendae: atque  subdit. ¶Vidua deligatur non minus sexaginta annorum: quae fuerit vnius viri vxor   in operibus bonis habens testimoniumt si liberos educauiti si hospitalis fuit, si sancto* rum pedes lauit si oppressis subuenit si omni bono operi inseruiuit. ¶ Quae tales habebant probitatis commendationes: dignae iudicabantur quibus de sacro subueniendum esset, etsi enim bona ecclesiae orphanis viduis et pauperibus dispensare bonum est: melius tamen vt sacra sacratioribus et meliora melioribus dispensentur. Quapropter sequenti praecepto iubet adolescentulas non esse hoc modo eligendas, etsi vnius fuerint vxores. nam plerunque euenit vt postq̃ bonis Christi deliciose fuerint vsae: nubere velint, in hoc ipso peccatum habentes: quod fidem viduitatis quam promiferant Christot per bigamiam violant. insuper cum sacris aluntur: discunt esse ociosae circuire domos  nugaces et curiosae et loquentes quae non decent. Ideo de eis quintum hoc dat praeceptum Paulus:dicens. ¶Iuniores autem viduas deuita. nam quando CHRISTO de,liciatae fuerint: nubere volunt habentes iudicium per primam fidem abiecerunt. simul  autem et ociosae discunt circuire domos: non solum autem ociosael sed et nugaces et curiosae lo quentes quae non decent. ¶Sunt qui haec duo vltima praecepta intelligunt: per facerdo\pend
\section*{COM }
\marginpar{[ p.5 ]}\pstart tibus ac presbyteris ad curandam rem familiarem liceat eligere viduas non iuniores q̃ sexagenarias. Sed nichil hici de re familiari presbyterorum loqui videtur Paulus: praeser tim cum tempore illolpresbyteri proprias vxores haberenti quae remiquae intus erati familiarem curarent. quid igitur opus erat de alienis mulieribus loqui: quae eorundem rem tractarent domesticam atque familiarem? De iis itaque loquitur: quibus ex sacro proui debatur. Quid est quando Christo deliciatae fuerint: nisi quando Christi deliciis vsae fuerint? quae sunt orare ecclesiae et rebus sacris interesse et insuper ex sacris educari, cum has tentauerint delicias: incipiunt plerumque iuniores viduae illas fastidire et nubere malunt. sicque  primam viduitatis fidem faciunt irritam: insuper et faciunt ecclesiae scandalum. Pro huiusmodi igitur iunioribus viduis: dat sequens praeceptum Paulus non tam iubenst q̃ indul.gens. Atque ait. ¶ Volo igitur iuniores nubere liberos procrearelrei familiari praeesse nul  iam occasionem dare aduersario conuicii gratia. nam iam quaedam abierunt retro post  sathanam. ¶ Quid est volo: nisi permitto. Praeceptum igitur hoc: permissorium est. Et melius est quod iuñiores nubant, q̃ fidem viduitatis promittentes ex sacris alimoniam petant: et postmodum fidem irritam facientes nubere velint. nam obiurgari solent plerumque huiusmodi: et diabolus ex obiurgatione occasionem sumens facit vt timentes huiusmodi irritae fidei conuicial non nubant et ad deteriora eas trahiti et secretas fornicationes. et sic auersae a Christo: eunt post sathanam et ruunt in perditionis barathrum. Et exemplo aliquarum quae sic perierunt admonet Timotheum: vt tales non recipiat sed eis indulgeat vt nubant. De iis viduis quae etsi iuniores sint et nubere nolint patres aut matres soceros aut nurus habentes: septimum tradit praeceptum Paulus dicens. ¶St  quis fidelis aut si qua fidelis habet viduas: ad sufficientiam suppeditet eis et non grauetur ecclesiai vt iis quae vere viduae funti sufficere possit. ¶ Haec de viduis: hactenus apostoli determinatio. Nunc de presbyterist id est fenioribus: qui verbo et exemplotid est sermone et opere bene praesunt sequens adiungit praeceptumi dicens. (Q ui bene  praesunt presbyteri: duplici honore digni habeantur maxime qui laborant in sermone , et doctrina. Dicit enim scriptura. Non frenabis bouem triturantem. Et dignus est ope, rarius mercede sua. ¶ Duplici honore: et spirituali et terreno. spirituali: eos reuerendo et venerando. terreno: eos nutriendo et omni corporali necessitate subleuando, quanto enim spiritualiores sunt: tanto quo ad se paucioribus terrenis contenti sunt, est enim magis deo simile: per se sufficientiat et paucissimis egere. qui autem terreni sunt: copiis sunt insaturabiles et semper crescit ardor habendi. Et maxime honorandi sunt: qui presuntt et in verbo fidei et in doctrina bene laborant. nam illi sunt architecti: alii autem vt iussa architecti exequentes. Hoc modo nos Christus hoc modo Christi imitatores apo stoli sermonis labore peruexere ad salutem. Et eos esse alendos authoritate sacrarum literarum firmat et veteris et nouae legis et vmbrael et veritatis. Nonum praeceptum de eisdem ponit: dicens. ¶ Aduersus presbyterum accusationem non admitte: praeterq̃ sub duobus aut tribus testibus. ¶Nam verisimile non est:vt qui ea aetate et in tanta gradus sublimitate positi sunti facile delinquant. et maxime in illa tempestate: in qua spiritus sanctus in cordibus eorum praesidebat. In omnibus enim qui spiritualiter praesunt: praesidet spiritus sandus. Nam illi: spiritu aguntur. Et age vt rem paruam magnae rei conferamus. Si quis praesit adolescentibus et ad mores et vitam instituat: aduersus eum non facile admitti debet accuatio. tanto minus aduersus seniorem qui praeest populo:nitens eum verbo et exemplo forma ,re ad pietatem. Decimum praeceptum subiungit: dicens. ¶ Peccantes: coram omnibus redarguel » vt et caeteri timore quatiantur. ¶ Materia subiecta vult id de pres byteris dictum. Qui sicut aduersus eos difficilius admittitur accusatio: ita conuicti et rei sceleris grauius debent redargui coram omniby aliis presbyteris (de presbyteris enim hic sermo est) vt et caeteri terreantur. nam quanto altiore in gradu sunt: tanto casus horrentior et magis timori esse debet. atque hac de causa: increpatio et deterritio maior. Qui vniuersaliter intelligunt: rem subiectam non attendunt. et id solum et maxime locum habere posset in publicanis et publice peccantibus. Ergo de sacerdotibus conuictis intelligatur: vt et consentaneum est. Sed et vndecimum quod maxime spectat ad episcopum: qui vt est eorum qui praesunt honorabilissimus, ita ad eum spectat labor sermonis et doctrinael et vitae exemplum: et cuiusque pro dignitate obsecratiot exhortatio approbatios aut redargutiol sed quemadmodum sancta mandant inslituta. Quapropter arcte mandat Timotheol et cum adiuratione: vt supradicta obseruet et nichil nisi quae praeiudicata praediffinita et preordinata sunt faciat motus aliorum aduocatione ac patrocinio , sic inqens.¶ Testor coram deo et domino IHESV CHRISTO et electis angelis: vt haec custo ,dias sine priori iudicio nichil faciens per aduocationem. ¶ Quid em sine priori iudicio:nisi quod praeiudicatum praediffinitum et praeordinatum a spiritu sancto et ducibus ecclesiae  \pend
\section*{X }
\section*{1Time }
\marginpar{[ p.210 ]}\pstart est? quorum vnus magnus fuit Paulus: in quo Christus et spiritus sanctus loquerentur. Quid est per aduocationem: nisi per humanam gratiam et fauorem? sic sacer interpretatur Athanasius. nituntur enim patronil aut aduocati humani quippiam accusantes aut defendentes: iudicis sibi conciliare gratiam. per quam si iudicat aliter q̃ a sanctis ordinationibus praesancitum est: ante deum et dominum Ihesum Christum et electos angelos  violationis sanctae adiurationis reus est. Oq̃ incumbit episcopis: recte iudicandi necessitas q̃ duras officiales eorum nunc exercent partes q̃ eis timenda est gehennae fornax et inextinguibilis clybanus. Horresco referens. Duodecimum praeceptum etiam ad episcopum pertinet: cuius est imponere manus et sacros fundere precatus et caetera exequi mysteria ad presbyteros consumandos. hi: dicti sunt reges sacrorum. neque enim decet regem aeternum filium aeterni regis: nisi per reges familiariter in altari tractari. et hi sunt per excellentiamtetsi non vulgari opinione reges: quos coronant angelil et quorum coronas angeli portantinon homines. de his verum est: et id. Et adorabunt eum omnes reges terrae. Ergo non temere non cito faciendi sunt presbyteri: ne scelestus aliquis tanta dignitate doneturi qua nulla in terris sublimiori nulla augustior. ne scelestus et impurus: proxime aeterni regis adhibeatur seruitio. alioqui qui ordinat: peccatorum ,eius reus tenebitur perinde atque scelerum particeps. Ideo subiungit Paulus. (I Nemi,ni cito manus imponas: neque communices peccatis alienis. ¶ Qui alienis vices suas tribuunt: operibus eorum communicant. vt qui iudex est: si vices suas alicui scelerato et iniusto committiti ipse pariter cum eo iniustitiarum quas substitutus commiserit, reus tenebitur. sic et episcopus sacram vicem augustissimil mysterii tractandilalii tradens: si spurcus et immundus fuerit indignitatum eius quoties spurce et immunde ad id mysterium tractandum accesserit pariter cum eo comparticeps reusque habebitur, sic et si ipse qui pastor est vicarium ponat ouibus suis qui vt lupus deuoret mutiler dissipet: peccatorum eius particeps habebitur et eiusdem reus tenebitur delicti. Item et si ipse qui iudex est officium alii committat qui iudicium peruertat: alienis communicabit peccatis, et reus erit peruersi iudicii. Et quia multa potest episcopus et variis modis potestatis aliis communicare: ideo multum vigilare debet et aciem mentis intendere ad id quod deo ple nus inquit apostolus neque communices peccatis alienis. nam haec sunt verba spiritus. Tertiumdecimum praeceptum communiter omnes respicit cum omnes regenerationis filii ad sanctitatem et puritatem vocati sint: et maxime episcopum qui summum locum inter eos tenet decet esse castum sid est purum mundum et sanctum. post quem: presbyteri. deinde , diaconi: vsque ad infimum plebis sanctae statum. Dicit ergo Paulus. ¶Teipsum: castum serua. (LVerum quia ad castitatem et puritatem seruandam: multum facit temperantia  et a cibis et potibus facientibus lasciuire et petulascere carnemtabstinentia. qualis esus multus carnium iusculentorum nimis delicatel opipare (et vt sic dixerim) Apiciane paratorum nimius spumante crathere et liberior vini haustus. Ideo subiungit Paulus ,ad Timotheum. ¶Noli adhuc bibere aquam: sed pauco vino vtere propter stomachum ,tuum et crebras infirmitates tuas. ¶Ergo vino et iis quae in sanis faciunt luxuriare car nem: vtendum est medicinae locot et cum natura ad sanitatem illa requiriti non ad lasciuiam. Certum enim esti sancti viri Timothei consuetudinem fuisse bibere aquam: alioqui non diceret apostolus noli adhuc bibere aquam. id enim insinuat noli adhuc. multum tamen vinumi prohibet: quia neque facit ad sanitatemi neque ad vitae castitatem atque puritatem. Et tales debent omnes esse: praesertim episcopi et omnes sacerdotes. qui si vino vtantur: pauco vti debent. Sed quid dicemus de multis qui immodice se proluunt: et plenis et turgentibus venis carnis puritatem perdunt pariter et spiritus ? et per omnimodam intemperantiam et exquisitos delicatosque cibos carnis petulantiam fouent ? et omnimoda gulae lenocinia sectantur vt videatur quidam intemperantiae daemon fomem ta gulael ventris renumque in popina miscuisse ad castitatem puritatemque perdendam? Quid enim dicemus de eiusmodi: nisi quod subiungit Paulus dicens? ¶Quorundam hominum peccata praemanifesta sunt: praecedentia ad iudicium. quosdam vero:etiam ? subsequuntur. Identidem et opera bona: quaedam praemanifesta sunt et quae aliter ha* bent: abscondi non possunt. ¶ Quid est quosdam vero etiam subsequuntur: nisi quorundam etiam peccata in fine deteguntur? Quid etiam et quae aliter habent abscondi non possunt: nisi quod bona quae ante manifestam non sunt in fine non latebunt. Et haec bona siue praemanifesta siue in fine cognita: ad iudicium vitae, identidem et mala (nisi expientur) siue i prin cipiossiue in fine manifestentur: in iudicium sunt mortis. Oculum enim nostrum: neque lux latet  neque tenebrae. stultum igitur est: putare siue bona siue mala quaecumque posse latere deum ma amum vniuersitatis rerum oculum. oculum vbique exillentem et omnibus praetentem, yidet enim  \pend
\section*{COM. }
\marginpar{[ p.5 ]}\pstart omnia: neque quicquam ipsum effugere potest. Studeamus igitur ita viuere: vt bona opera nostra manifesta semper sint ante eum etsi mundum lateant: quae etiam in fine latere non poterunt mundum. vt ei placeamus: qui cum sit summe bonus amator est bonorum. et ei seruiamus: quem boni omnes sine fine laudabunt. Cui et ex hoc nuncisit omnis laus  decus et honor:in omnia saecula saeculorum.  \pend\pstart  (EXAMINATIO nomnullorum circa literam. ¶ Vulgata aeditio. Viduas honora:quae vere viduae sunt. ¶ Vere: aduerbium est non nomen. Paulus. χή ρασ τίμα, τὰς ουτωσ χήρας. Et sic vere: vbicunque in hoc capiteliungitur vocabulo viduael sumitur. , ¶ Vulgata aeditio. Discat primum domum suam regere et mutuam vicem rependere paren,tibus. ¶Dicendum: discant pluratiue. Filiis enim et nepotibus qui sufficienti et adulta sunt aetate: praeceptum dat Paulus non viduae solatio viri destitutae, qd̃ insinuat particula: et mutuam vicem reddere parentibus. id eni: ad filios et nepotes pertinet. Et quod vetus ĩterpres dicit regere: magis habetur pie religioseque colere. Eusebia enim (quae pietas et pia cura est) ad deum ad parentes atque ad patriam dicitur. Paulus. μαμθαρέ τωσαν τὸν IAιoμ ὁικομ ἐυσεφςῖμͅ. καὶ ἀμυοιθας ἀττολιλόpαι τοῖσ ωρογόμοισ. ¶Vulgata aedi* tio. Quae autem vere vidua est et desolata:speret in deum et instet obsecrationibus. ¶Id desolata: solitarie manentem significat. Et dicendum potius sperauit et immaneti slue insistit orationibus. Quia enim sperauit: sola manet libenter et orationibus iugis insistit. Etiam si diceretur sperat in deum et imanet orationibus propter naturam copulatiuae particulae quae id apud latinos vtplurimum idem tempus requirit:neque forte officeret intelligentiaei neque ser monitetsi apud Graecos sperauit praeteritum sit aptusque sit sermo. Paulus. iiλὲ ὀprωσ xήρα, καὶ μεμορομeρη, ἠλπικερ εὐωι τὸμ οεὸν, καὶ πpοσuepeι, rαῖς Aς ,uσεσι. ¶ Vulgata aeditio. Si quis autem suorum et maxime domesticorum curam non habet: ,fidem negauit. ¶ Potius dicendum: suarum et domesticarum foeminino genere. nam de viduis loquitur Pauls. Eniuero ꝗ curam non gerunt aut matrum aut filiarum, aut consanguinearum viduarumi sed vtcumque volunti sinunt licentiose viuere et delitiis ac voluptatibus defluere quae alioqui prouidentia suorum probae et honestae viduae fuissent: peiores sunt infidelibus qui nunq̃ fidem agnouerunt et durius perferent iudicium. Paulus. eιλέτισ τῶν IAIωp, Rdl̀ μάλιςα τῶμ ὀιπείωμ ὀυ πρόμοετ͂, τὴ̓μ πίσιμ uppuται. Quia enim genus apud Paulum ambiguum est (nam commune) masculinum coepit interpres apud nos minime dubium, sed quod minus pro subiecta materia videbatur. Saltem vt ancipitem inuenerat sermonent ancipitem interpretando reliquisset vt et nos fecimus. qd̃ discretio legentis: in id quod oportuerat discreuisset. ¶ Vulgata aeditio. Cum enim luxuriatae fuerint: in CHRISTO nubere volunt. ¶Dicendum esse potius videtur: deliciatae Christo. Et cum dixiti iuniores autem viduas deuita: non tam aetatem respexisse crediderim q̃ morum mobilitatem et inconstantiam. Nam si moribus iuuenes funt et inconstantes: volunt vt vere viduae quodam temporelsolae maneres orationibus et obsecrationibus insistere. sed vbi parumper hoc viuendi modo Christo deliciatae sunt: incipiunt fastidire et nubere volunt. Verum ob mo rum inconstantiam quia moribus iuuenes sunt: si delicias Christi eligunt quae sunt in orationibus et obsecrationibus et abstinentiis degere statim fastiditae desistunt et ociosae manent et discunt circuire domos fieri verbosae et loqui quae non decent. Tales sunt deuitandae neque etiam a pontificibus (etsi id peterent aut earum parentes) seruitio Chri sti in claustris mancipandae. Talibus profecto: ex bonis ecclesiae alimonia non debetur. Sed parum nunc praeceptum Pauli obseruatur: et parum quae tales sunt deuitantur. Et in superiori sermone mutauit vetus interpres sententiam: aut insulsami et indignam fecit locur tionem. Nam cum vocabulum Christot in verbum praecedens apud Paulum referatur: nemo nisi indigne et insulse dicet luxuriatae Christoi siue in Christo. Quid enim luxuriatae: nisi for nicatae ( at nullus dicet fornicatae Christo aut in Christo. Qua de causal mutata sentential praepositionem addunt Christo et in sequens verbum nubere reuocantsatque sententiam mu tant. Paulus. 8rαp γαρ καταςρμμιάσωσι τοῦ χριςοὺ, γαμεῖῦͅ θε λουσιυ. ¶T Vul,gata aeditio. Volo autem iuniores nubere. ¶ Dicendum: volo igitur. Concludit enim hic quid de huiusmodi iunioribus non tam aetateq̃ moribus faciendum sit: quas superius ad delicias Christi iussit esse deuitandas. Paulus. όυλομαι ο͂υμ μεωτέρασ γαμεῖν¶ Vulgata aeditio. Si quis fidelis habet viduas: subministret illis. (T Dicendum potius: si quis fidelis aut si qua fidelis. Nam fidelis prius in masculino et consequenter foemininum ei subiungitur. Paulus. 21Tισ ωιSὸσ ηι πιςὴ ἐχςι χἠρασ, ἐπαpRειτω ἀυταῖσ. Verum si diceremus si quis fidelis aut fidelis habet viduas: secundo loco genus non intelligeretur( vt in Graeco intelligitur iccirco adiecimus si qua. Sie enim genus manifeste cum primo tum secundo loco per articulum in graeco intelligitur sine tali adie\pend
\section*{x }
\section*{1 Timo. }
\marginpar{[ p.211 ]}\pstart cione: acsi diceremus si quis fidus aut fida viduas habet, sed fidus: non tam accommodate exprimit vocabulum vt fidelis. ¶ Vulgata aeditio. Vt non grauetur ecclesia. ¶Legendum per copulandi notam: et non grauetur ecclesia. Paulus. Και μὴι δοαρέ ισθων ηι ἐκ*λησία. ¶LVulgata aeditio. Testor coram deo et CHRISTO IHESV. (J. Dicendum: et domino Ihesu Christo. Deest domino et ordo sermonis immutatus. Paulus. AιαμαpΤύρομαι ἐρώ τrιoμ τοὺ οεοὺͅ, καὶ Κυριoυ l̓ugoὺ yριςoῦ. ¶ Vulgata aeditio, Vt , haec custodias sine praeiudicio nichil faciens. ¶ Particula sine praeiudicio: in particulam sequentem referri debet. Et sine praeiudicio: non id significat quod vulgo intelligi solet sine detrimento aut dispendiol sed id significat sine praeuiol priore praecedenteue iudicio. designat enim praeiudicium: praeuium praecedensue iudicium. Quam intelligentiam: in expositione huius locil haud inscite beatus sectatus est Ambrosius. et id: ratio vocabuli exposcit. Paul, i̓́μα ταῦτα φιλάςησ, χωρὶσ τροκρίματος μηλὲμ ποιῶμ. » ¶ Vulgata aeditio. In alteram partem declinando. (I Potius dicendum: per aduocationem. Nam ωρόσ λλησις non declinationem: sed aduocationem significat, quod perinde est: ac per patrocinium aliorum qui aut aliqua accusant aut aliqua probant. Clisis enim quando declinationem significat: scribitur per iotatid est per iinonam literam. quando vero vocationemn: per itatid est ul septimam literam graecam. hic autem per ita septimam literam scribitur. Vicinitas tamen vocabulorum: forte fecit veteri interpreti allucinationem vt longa circulatione vteretur et paraphrasi non interpretatione. Paulus. κατὰ ωροσΚAuGιp. ¶ Vulgata aeditio. Quorundam hominum peccata manifesta sunt. ¶ Po tius dicendum: praemanifesta sunt. Paulus. τιμῶμ ἀμορό ωωμ ἀι ἀμαρτ ἰ αι πρόν Au λoi ἐισι. ¶Vulgata aeditio. Similiter et facta bona: inanifesta sunt. ¶I Et hic quoque  dicendum: praemanifesta sunt. Paulus. ὡ̓σάυτως καὶ τὰ καλα ἐρτα, πρόAηλά 251.  \pend
\endnumbering\beginnumbering\section{CAP. VI.}\pstart \huge\textbf{D}\normalsize E viduis presbyteris et episcopis: posuit in praecedenti capite praecepta. Nunc de seruis: hoc ponit praeceptum Paulus dicens. ¶Q uicunque sub iugo sunt serui: proprios dominos omni honore dignos arbitrentur enim nomen domini et doctrina blasphemetur. ¶ De iis seruis hoc praeceptum intelligitur: qui sub iugo seruitutis erant dominorum infidelium. Quis enim dubitat si domini nossent seruos effectos Christianos et tunc minus seruire? et tunc minus se ab eis in honore haberi: quin nomen et doctrinam Christi blasphemarent cul credentes serui et cuius doctrina imbuti tam malii rebelles superbisac contemptores ef, fecti essent? negligentia igitur eorum et contemptus: in contumeliam Christil et doctrinae eius reciderent: nisi (vt par est) seruirent et dominos suos omni honore dignos du o cerent. et hoc propter Christum: qui est dominus dominorum. Aliud praeceptum dat de  iis seruis: qui dominos habent fideles dicens. ¶ ¶ Qui vero fideles habent dominos: » non contemnant quia fratres sunt sed magis seruiant quia fideles et dilectit qui beneficium percipiunt. ¶ Fideles serui fideles habentes dominos: non debent vt seruil conditionem attendere spiritus. nam secundum conditionem spiritus qua Christiant sunt: non seruil sed fratres. et domini: secundum conditionem spiritus non domini eorum sed fratres sunt. sed serui conditionem carnis attendere debent: qua et ipsi reuera sunt seruil et domini eorum domini. Et nichil refert: an empti fuerint aut vernaculi aut lege humana serui. nam hanc conditionem voco carnis: quae ex spiritu non esti sed aliunde. et conditio spiritus hanc non tollit: quia peccatum non est et sine peccato esse potest. et seruus: sine peccato seruit. et dums: sine peccato dominatur. neque seruire neque dominari vt huiusmodi: peccatum est. Quid igitur attem dere debent huiusmodi serui ? Conditionem profecto carnis: in seruiendo et seipsos omnimodo humiliando ac subiiciendo. et conditionem spiritus: ad diligendum pariter et seruitium diligentius exequendum. Quo modo enim eos non amabunt: quos secundum hanc conditionem cognoscent esse fratres ? quo modo cum conditionem carnis et fortunae attendunti aut fortasse secreti cuiusdam diuini iudicii: non promptius seruient iis quos dilectos esse cognoscunt? Sed cui dilectos? Christo et regi aeternol et vniuersorum domino. Multo sane magis est eis occasio seruiem di q̃ infidelibus dominis: qui etsi in se potestatis et dominii ferant signaculum (omnis enim potestas a domino deo est) non tamen dilecti sunt a domino. Haec doctrina: spiritus est. Quapropter mandatum dat Timotheotimo per Timotheum omni praecont der: vt ita ,doceat atque ait. ¶ Haec doce: et exhortare. Si quis secus doceti et non accedit sanis sermonibus qui sunt domini nostri IHESV CHRISTI et doctrinae quae secundum pie,tate est: ianiter se se iactauit nichil sciens. sed languens circa quaestiones et pugnas verborumex quibus oritur muidia contentio blasphemie suspitiones malaes conflictationes  \pend
\section*{COM }
\marginpar{[ p.6 ]}\pstart . hominum mente corruptorum et veritate destitutorum existimantium quaestum esse pie* tatem. (¶Praecepta quae posuit Christiloquus et spiritiloquus Paulus et quae iubet esse do cenda et hominibus suadenda: non suos sermones esse diciti sed spiritus sermones et sermones domini nostri Ihesu Christi. et qui aliter docent: inaniter se de doctrina iactant, quia nichil sciunt: etsi se scire putant. Et vt eos cognoscamus: manifesta dat nobis signa. nam ii sunt qui circa questiones languent et logomachias id est argutiarum verborumo pugnas: ex quibus inuidialcontentio blasphemiae ebulliunt. Nam quod vnus de diuinis asseuerat: alter continue pernegat. quare si vnus bene dixit: alter blasphemat. Suspi tiones malae. nam vicissim contendentes:se male vicissim de fide plerumque sentire suspicantur. Conflictationes hominum mente corruptorum. nam vt ii qui appetitu corrupti vnguibus et manibus aut etiam laetiferis armis in se inuicem conuolant et seuiunt: ita qui mente corrupti suntltam vana et inutili verborum pugna ad inanem gloriami aut quaestum confligunt. et quo modo corrupti mente non essent: qui veritate destituti sunti existimantes quaestum esse pietatem? Veritas enim: est mentis vita. Quid est quod doctrinam spiritus: aut prorsus abiiciunt aut carnalitersaut infra carnalem etiam dignitatem trahunt? Vtinam tempora nostra nullos tales pariant: et praesens aetas talibus careat monstris. Oi si erunt: quid nam faciendum longe etiam per spiritum videns: sequenti documen. to apostolus nos edocer dicens. ¶ Procul esto a talibus. Est certe quaestus ingens: pietas cum sufficientia. nam nichil intulimus in hunc mundum: manifestum pro neque quicquam ,auferre possumus. habentes autem alimenta et integumenta:eis contenti simus. ¶ Pugnas verborum inanes reliquamus: quae proculdubio non pariunt nisi verba. et mera verba sunt: meraque verbositas. alioqui inanes non essent. et doctrinas sectemur pietatis: quae sunt doctrinae spiritus. Ex illis quae verborum pugnae ac conflictus plenae sunt: inanem gloriam et auaritiae quaestum non quaeramus. sed doctrinam sectantes pietatis: pietatem consequemurt ingentem reuera quaestum et per se sufficientemi aut saltem minimil non sui quidem sed ratione carnis egentem. propter quam alimento eget non multot sed moderato: et vestimento quo tecta caro defendatur a nociuis, non quod superbiat pauonis more: quae volucris in se picturam gerit angeli in se conuersi e coelo ruentis. Sed quid i externis bonis gloriaris:in vestibusl in diuitiis etiam si ex diuina benedictione et propriis rebus sine auaritia vt Iobi vt Abrahans vt Isaacivt Iacobl tibi assint? Cum cogitas te nichil tecum in hunc mundum ex illis omnibus intulissel neque ex illis quicquam tecum auferes in alterum? parum tibi gloriandum: nisi rationis oculo captus sis qui longe etiam inferior est ocu lo spiritus. Qz si diuitiis non abundas neque habes vnde ex dei benedictione et prouidentia accipias: quid illis inhias quas sine auaritia habere non possis Quo modo sine auaritia ha bere putas ?cum ille appetitus tuus rationis corruptae auaritia sit ? Esto igitur i studio pieassunt diues fieri: et audias apostolum id mandantem, atque dicentem. ¶Qui enim volunt diuites fieri: incidunt i tentationem et laqueum et multas irrationabiles et noxias concupiscentias: quae demergunt homies in interitum et perditionem. radix enim omnium malorum: pecuniae cupiditas * est. quam quidam appetentes aberrauerunt a fide: et seipsos iplicuerunt doloribus multis. ¶Oz radix omnium malorum sit philargyriai id est argenti pecuniaeque amor atque cupedia: clarum est. Nam auertit ab amore dei: et amorem qui sua natura coelestis est deflecit detinetque in terram non aliter q̃ sepum coelestem flammam quasi auaram inferioris pabuli detinet ad ima, quae si emicare sursum nititur: statim extiguitur. ita et auari mens: si sursum tendere vult non sentit amorem. imersus enim rebus terrenis: in altis viuere non potest. Extinguit igit pecunie amort charitatem dei: quod maximum et capitale malum est. extinguit charitatem ad proximum: vt cuius amplioribo fortunis iuidet. aut bonis eius si mediocris fortunae est inhiat: aut si omnino pauper smolestus est enim de rebus eius quicquam petat aut absumat. Ne lucrum non faciant substam tiae terrenae: sabbata violant et nichiliducunt lucrum perdere vitae aeternae. Si parentes viuunt? si opulenti sunt: vita eorum eis tedium affert. et mortem et patris et matris optant: vt diuitiis eorum potiantur augeantque ifoelicis auaritiae suae cumulum. Ecce quo modo parentes ĩhonorant. Sed quot iam: propter pecuniae amorem thori sacramentum violarunt tam viri q̃ mulieres? Ar bitror apud iferos illorum miserrimum esse eiulatum. Vnde furtal vnde aliorum opum concupiscem tiae: nisi ex insano pecuniae amore? Quotusquisque nunc magistratus et honores ambit: vt queat diues fieril neque santis vnq̃ marsupia nummis abundant? Non est species malil quam haec virosa radix quam diabolus fixit non proferat: vt admodum certum sit quod spiritus locutus est in Paulo radicem omnium malorum esse pecuniae cupiditatem. pecunia id malum non esti radix illa non est: sed ipsa circa pecuniam cupiditas quam summopere fugere oportet. Neque satis est eam fugere: nisi etiam sequamur quae dei preco apostolus hic iubet esse sequenda, cum ait.  \pend
\section*{X }
\section*{1 Timo. }
\marginpar{[ p.212 ]}\pstart ¶Tu autem o homo dei: haec fuge. Prosequere: iustitiam pietatem fidem dilectionem  *patientiam mansuetudinem. Certa bonum certamen fidei. Appraehende vitam aeternam: ad quam et vocatus es et confessus es bonam confessionem sin conspectu multorum testium. ¶Certamen Timothei erat contra falsitatem erat contra infidelitatem: quam vincendo et extirpando vita premiabatur aeterna. Hec enim erat bona dei voluntas. Confessionem bonam confessus est Timotheus: in conspectu multorum testium. nam publice fatebatur dominum nostrum Ihesumtomnium saluatorem natum de virginetin terris visum pro peccatis nostris mortuumspro gloria nostra a mortuis surrexisse et reliqua om nia quae sanctissimae fidei sunt consentanea. Nunc arctissime praecipit Paulus eidem Timotheo: vt quod sibi mandatum est inuiolabiliter seruet. Sic quodque dei mandatum inuiolabiliter custodire debet quisque episcopus: vt qui rationem eius reit in die aduentus domini sit redditurus. quem ille qui solus beatus est potens est inuisibilis ( et immorta lis rex et dominus omnium suo tempore visibiliter manifestabit. Ait ergo. ¶ Praecipio ti*bi ante deum viuificantem omnia et CHRISTVM IHES VMIqui testatus est sub Pon *tio Pilato bonam confessionem: vt serues praeceptum immaculatum sine repraehensionel * vsque ; ad apparitionem domini nostri IHESV CHRISTI. quam propriis temporibus: *ostendet beatus et solus potens rex regnantium et dominus dominantium qui solus *habet immortalitateml lucem habitans inaccessibilem. quem nemo vidit hominuml neque videre  potest. Cul: honor et potestas aeterna. Amen. ¶ Vt deus omnibus dat esse: ita omnibus qui viuunt dat viuere. Nam ipse solus est omnium entitas et omnium vita. Sed age: quam bonam confessionem testatus est dominus noster saluator omnium et viuificatori sub Pontio Pi lato? Rex sum ego. ego in hoc natus sum et ad hoc veni in mundum: vt testimonium perhibeam veritati, omnis qui est ex veritate: audit vocem meam. Vere rex noster em et omnium: qui in hunc mundum venit et nos omnem veritatem quae aestimanda est docuit et om nes qui sequuntur eum. quia audierunt vocem eius: et obediunt praeceptis eius. Ipse est rex regnantium et dominus dominantium: vt in sacra visione vidit dilectus dei discipulus. Vestitus (inqt) erat veste aspersa sanguine: et vocabatur nomem eius verbum dei. Et subdit. Et habet in vestimento et in femore suo scriptum: rex regum et dominus dominantium. Praeterea: haecl et id beatus solus potens: communia habet cum patre et spiritu sanctoiet theologicam cataphasim respiciunt. Quae autem ex sequentibus intelliguntur solus immortalis inaccessibilis inuisibilis impossibilisque videri: ad sacram de eo pertinent ignorantiam sacrumque non solum in terra sed et in coelo silentium et ad theologiam quam sacre philosophantium dux vnus Dionysius appellat apophaticam. Et haec vniuersa: absolute et incontracte deo conueniunt. caeteris autem si qua conueniunt vt beatitudo potentia maiestas dominatio imortalitas inuisibilitas: participatiue et contracte conuent unt. at sic participantia cum absolutis et contracta cum incontractis consentiunt: vt finita cum infinitis. Est enim visus in corpore: intellectus contractus. et sensualis appetitus: contracta voluntas. at q̃ corporeus visus non est intellectus et sensualis appetitus non est voluntas: tam angeli et homines si ad deum in absolutione conferas simo infinito plus neque beati neque potentes neque sunti neque esse possunt. sunt tamen et possunt: suo modo. Nam cuiusque modus: contractio est. ergo deus absolute talis sine modo supraque  omnem modum. Et vt varie intellectus in variis contrahituri nam aliter in visione ali ter in imaginatione: ita beatitudol et potentia) et reliqua id genus quae non nisi contracte creaturis tribuuntur varie in angelo et varie in homine contrahuntur. Sed (vt par est) haec omnia incontracta: soli deo tribuit apostolus. Verum iis dimissis quae ad altissimum philosophandi de deo modum spectant: reuertitur ad dandum Timotheo pro diuitibus huius saeculi praeceptum. Nam et ipsi saluari possunt: si oculum spemque conuertunt. Terrena attendentes superbiunt: tanto magis humilientur coelestia attendentes quae ipsi perdunt. Opes respiciunt quia bonae sunt: deum respiciant quia summum bonum et superans omne bonum. Opes amanti quia delectabiles sunt: multo magis ament deum quia summum est delectabile. Confidunt in opibus quia ab indigentiis in hoc mundo liberant: se conuertant et confidant in deo quia ab omnibus indigentiis non modo in hoc mundo  sed et i altero vere liberat ac solus liberare porest. solus dat veram sufficientiam: qua habital nullum vltra torquet cruciatque animum habendi desyderium. bona eius sunt constantissima cer tissima enim ireptibilia : vitam non parui temporis sed aeternam praebentia. Hoc igitur pro eiusmodi : sanctum et salutare (si audiunt) ponit praeceptum apostolus dicens. ¶Diuitibus ꝗ in presenti ,sunt saeculo praecipe: vt non alte sentianti neque sperent in incertitudine diuitiarum sed in  deo viuo qui praebet nobis omnia opulente ad delectabilem vsum vt bene agant?  vt diuites sint in bonis operibus vt facile contribuant. communicatiui thesaurizantes  \pend
\section*{COM. }
\marginpar{[ p.6 ]}\pstart ,sibiipsis fundamentum bonum in futurum: vt aeternan vitam appraehendant. ¶Hoc praeceptum multa pariter complectitur. Primum: vt diuites non alte non superbe sapiant. nam diuitias huius saeculi comitari solet superbia: quae nisi per humilitatem cordis repressa fuerit ipsi fructum bonum facere non possunt neque venire ad vitam. est igitur illa de se alte sentiendi opinio: ante omnia extirpanda. Secundum: vt non sperent in incertitudine diuitiarum suarum. incerti sunt quantum sibi duraturae sunt: immo quantum ipsi suis diuitiis duraturi. quo modo igitur i rebus vtrinque tam incertis: sperare possunt? sed in deo viuo diuitiarum omnium largitore ad vere delectabilem vsum: confidant qui incertum non est bonum qui omnibus semper durat et cui omnes semper durabunt aut in fruitione aeternorum bonorum aut in gehemna tormentorum. Tertium: vt ex diuitiis suis bene agant. non fecit deus soleml vt sit ociosus: ita neque diuitias vt sint ociosae. sed omnia illat fecit propter vsum. Diues homo qui bonis suis non vtitur: accensum lumen claudit in archar accenso etiam luminel semper noctem habens. Liberalis contra: semper lumen habet in apertoshaud secus q̃ si in sole esset numq̃ noctem sentiens. nescis enim q qui in coelo habi. tat noctium et dierum vices non habet: sed semper est in die? Quartum: vt bonis operibus sint diuites. Nam aurum argentum census? et possessiones non sunt verae diuitiae: sed bona opera maximel quae diuinas gratias comitantur. ac deinde virtutes nostrae industriae earumque opera. Q uintum: faciles sint et prompti contribuendol et bene contribuendo. Nam qui aegre et difficulter id faciunt: adhuc auaritiae veterno languent et nondum amputata est pecuniae cupiditas quae est radix omnium malorum. Sextum: diuites sint commu nicatiui. sic enim melius intelligunt: quibus bonorum suorum beneficia digne sint impartituri. Et haec facientes: vere diuites erunt recondentes sibi fundamentum bonum vitae futuraetin aeterna aedificatione. Vltimum praeceptum Timotheo apponit. Vt talentum sibi creditum gratiam sibi depositam custodiat: et haereses euertat et destruat. et i hoc valefacit ,sancta vt solet et pia valedictione. Ait itaque . ¶O Timothee depositum custodi: euertens  prophanas vocum vanitates et oppositiones gnoseos falso nominatae. quam quidam pro,mittentes: a fide aberrauerunt. Gratia tecum. Amem. Ad Timotheum primal scripta Laodiciae. ¶ Episcoporum est: fideles instruere auersos conuertere et ex ifidelibus constituere fideles  haereses euertere et omnem vanam eruditionem funditus extirpare sanctorum apostolorum qui haec omnia sublimiter et valide faciebant (remota licet) imitatione. qd̃ et de Gnosticorum vanitate: hic apostolus Timotheo iubet faciendum. Gnostici enim authore Irenaeo: a Simone mago profluxerunt et a gnosit id est cognitione seu scientiai sed falso nomine (vt ait apostolus) dictal Gnostici etiam falso nomine dicti sunt. vocant enim se cognitiuos: et nichil cognoscunt. scientiam fidei promittunt: et a fide aberrant. Et ho rum vocabula yt et Simonis magi: cenophoniaei id est vanae voces sunt et a re signifit cata vacuael multo magis q̃ Chimeral et Hippocentauri. Simon magus figmentum ponebat virtutem quandam conditore mundi superiorem quam deum appellabat: estque  deus Simonis magi nomem inane. et de deo vero dei nomen negat: vt figmento tribuat. Sic figmenta sunt apud Gnosticos Barbelol vt virtus masculai et ennaea vt foemineat et tota proles Prognosis aphtharsiai et Zoel et phos: et nus et logos et thelema. Quid Barbelo sonct supremae virtutis eorum figmentum: incertum. caetera latine dici possunt? intelligentia praescientia) incorruptibilitas vitat lux mens ratio voluntas: pulchra profecto vocabula sed (vt ipsi vtuntur) vana prorsus( et fictitia. de veris enim negant: vt figmentis et quae nichil funt et nichil esse possunt tribuant. Coniugia etiam figmenta iunt daemonum: ennaeae et logit aphtharsiae et photos quem Christum appellant Zoes et thelematos nus et prognoseos, id est coniugia intelligentiae et rationist incorruptibilitatis et lucis vitae et voluntatis mentis et praescientiae. Et aeones id est saecula omnia ipsorum: diabolica sunt figmenta. nomina quidem ma gnificai et curiosa: res vero nulla iuxta eorum intelligentiam respondens neque respondere potens etsi pu tent aliquid respondere. sic eos figmentis illusit diabolus. Haec magnificorum nominum glumal non solum Gnosticos et eos qui ab eis seducti sunt decepit: sed et Platonicos multos de vnosideal et anima mundi. quae ad Platonicorum intelli gentiam (vt verbo Pauli vtar) cenophoniae sunt et voces vacuae rebus (pace eorum dixissem) sed malo pacem Christi. quam illis omnibus bene opto: et sancti spiritus sequi doctrinam. Sed diuinus ille Socrates asseruit. Diuinum dixerim: si quem placet dici diuinumi qui edoctus sit a daemonio. At quid nam aliud daemon doceret de diuinis: nisi ea quae sunt maxime ardua et magnificatad figmenta deducere? odit enim veram diuinitatem: et gaudet fictam et ementitam inducere. et magnifico verborum apparatu: nobiliores ludificare sensus. At multilinguis hydrus Iibilat: Augustinum celebrasse ideame  \pend
\section*{x }
\section*{1 Timot }
\marginpar{[ p.213 ]}\pstart Esto. sed plane (vt vir sanctus) ad catholicam non ad Platonicorum intelligentiam ideae nomen Christo filio dei tribuit. quem Porphyrius Plotinus Iamblicus Proculus et caetera Platonicorum turba non negant modo: sed persequuntur eiusque aduer santur doctrinam. quorum libri a recentioribus Platonicis: etiam nunc ad astra feruntur. Et vt Augustinus: sic etiam celebro intelligentiam incorruptibilitatem vitam lucem mentem rationemi diuinamque voluntatem: non ad Gnosticorumi sed ad Christianami sanctam piamque  intelligentiam, non enim in verbis peccatum: sed in ĩtelligentia verborum mentisque adaptatione. Sinamus igitur hos et illos et omnes vaniloquos valere: et omnes contra fidem obices euertamus eos a nobis protinus repellendol vt iniectos sinibus angues. et sequamur Christum regem nostrum vera cognitione et scientia: eum ex studio sacrorum eloquiorum celebrantes et nunq̃ vana curiositate aut mentis praesumptione a semita spiritus sancti discedentes vnum deum in substantial et trinum in hypostaticis beatitudinibus confitentes. Quem decet omnium creaturarum laus perennis immortale decus verendum imperium regnum omnium seculorum et gloria incompraehensibilis: in omnia saeculorum saecula. Amen.  \pend\pstart *¶EXAMINATIO nonnullorum circa literam. ¶ Vulgata aeditio. Quia beneficii participes sunt. ¶ Legendum:qui beneficii. Paulus. 6̔ι τὴσ ἐυεργεσίας ἀντιλαμφαυόο . ρόμεpοι. ¶ Vulgata aeditio. Superbus nichil sciens sed languens circa questiones et pu» gnas verborum. ¶ Quod vetus interpres dicit superbus: verbum est praeteriti temporis apud Paulum quod superbire et inaniter se iactare significat, qui enim aliter docebant q̃ apostoli tradiderant: superbiebant docentes audentes ac praesumentes quae docenda non erant. et nequicquam vane iactabant: se ea scire quae nesciebant. O quot et hodie sunt qui relictis apostolorum doctrinis circa diuina superbiunt: se vanis in rebus iactantes circa questiones pugnasque verborum elanguentes. qui de numero eorum sunt: quos spiritus domini cuidam religioso viros fratri Roberto in spirituali visione ostendit lapi dem magnum et oblongum rodentes. ex quo duorum capita serpentum exertis micantium linguis: prominebant. quam vt legentes ab huiusmodi inanibus friuolis et infructuosis resipiscerent studiis: veram esse puto. Optimum sane est:in sacrarum literarum doctrina se continere et earum intelligentiae non humana praesumptione quaesitae sed diuini spiritus illapsione insistere aut Petri in Clemente sequi sententiam. Diligenter (inquit) obseruandum est: vt lex dei cum legitur mon secundum proprii ingenii intelligentiam legatur. sunt enim multa verba in scripturis diuinis: quae possunt trahi ad eum sensum quem sibi vnusquisque sponte praesumpsit. quod fieri non oportet. non enim sensum quem extrinsecus attuleris alienum et extraneum debes quaerere quomodo ex scripturarum authoritate confirmes: sed ex ipsis scripturis sensum capere veritatis. Et ideo oportet ab eo intelligentiam discere scripturarum: qui eam a maioribus secundum veritatem sibi tra ditam seruant, vt ipse possit ea quae recte suscipit: competenter asserere. Cum enim ex di uinis scripturis integram quis susceperit et firmam regulam veritatis: absurdum non erit si aliquid etiam ex eruditione communitac liberalibus studiis (quae forte in pueritia attigit ) ad assertionem veri dogmatis conferat. ita tamen vt vbi vera didicit: falsa et simulata declinet. Haec Petrus. Sunt praeterea sacrae literae delicatissimus panis et sacra spiritualis intelligentia optimum coelesteque vinum: et vtraque vitalis et saluberrimus ani mae nostrae cibus, quem qui dimittunt ob inanes verborum pugnas: anguiparum proculdubio ( nichil vitaliter proficientes) erodunt faxum. similes iss quos Prudentius vir deo dignus in Apotheosi exccratur dicens. Soluunt ligantque questionum vincula: Per syllogismos plectiles. Vae captiosis sycophantarum strophis. Vae versipelli astutiae.  \pend\pstart Sed si Paulus Christi dei aeternaeque ; sapientiae apostolus ssi dei spiritus in prophetico viro / si Petrus ecclesiae columen si Prudentius immo nullus prudens audiuntur: quomo do nostrum valebit consilium? Nisi certe Christus erigat:ingenia prostrata sunt in imas se sua sponte deuoluentia latebras inanesque nenias. Et scimus qui illa scabie pruriunt: eis nostra monita non placeret non aliter q̃ infirmis medicorum antidotassalutaresque phar matias. Sed non ob idipsum: desinit medie admonere et si potest etiam sanare. Nos nul li inuidemus nulli malum volumus: sed si possumus magnum et salutare bonum. Pau lus. τετ ύφωται μηλὲρἐτιςάμερος, ἀλλὰ υοθῶμ περὶ 3ητήσίσ και Aoyo,μαχίασ. ¶ Vulgata aeditio. Ex quibus oriuntur inuidiaescontentiones. ¶ Ex Paulo singulariter dicitur est fit aggignitur siue ortur iuidia lis fiue contentio. Paulus.es op  \pend
\section*{COM }
\marginpar{[ p.6 ]}\pstart  γίρεται φοονοσ, ἐρισ. ¶Vulgata aeditio. Existimantium quaestum esse pietatem, est autem quaestus magnus. ¶ Continuo post pietatem: deest enim particula integrat absistito siue procul esto a talibus. Paulus. poμιςόρτωμ ττορισμὸμ ε̃͂ιραι τη̃μ ἐυgέpειαν, ἀφίςατο ἀπὸ τῶμ τοιούτωp, ἐςι λὲ τορισμοσ μέγασ. ¶Vulgata aeditio. Inci ,dunt in tentationem et laqueum diaboli et desyderia multa. ( Diaboli: abundat etsi intelli gentiae nichil officit. Paulus. ἐμιτί πουσιμ ἐισ πειρασμὸν, καὶ πάγι λαν καὶ ἐττια σουμίας πολλάς. ¶ Vulgata aeditio. Inutilia et nociua. ¶ Potius dicendum: amentia/  siue irrationabilia et nociua. Paulus. αuou τουσ Και φλαθεράσ. ¶ Vulgata aeditio. » Radix enim omnium malorum em cupiditas. ¶ Cupiditas generalius exprimit q̃ Paulo dicit. Nam vocabulum Pauli specialius designat argenti nummil pecunieue cupiditatem. qua maxime fer uent qui volunt diuites fieri: et quam quidem foedam animi perturbationent nostri enim ferme vo cabulo auaritiam nominant. Paulus. pi τα γἀρ τάυτωυ τῶμ Κακῶμ ἐςίμ η φιλαρuε ,pia. ¶ Vulgata editio. Vsque in aduentum domini nostri IHESVCHRISTI. ¶TPotius: apparitionem q̃ aduentum hic Paulus habet. Paulus. μexρι τησ ἐπιφαμείας τοῦ κυρίou » ἡμῶμ ἰμ̓σοὺ χρισοῦ. ¶ Vulgata aeditio. Rex regum et dominus dominantium. ¶ Potuit vets iterpres dicere rex regnantium vt dixit dums dominantium. Cur enim non aequo iure primo loco purticipium p participium iterpretatus entvt secundo subinde fecit loco? Paulus. Pασι λEUσ TOU  Φασι λευόυτων, Και kύριοσ τῶμ kυρι eυópTOU. ¶ Vulgata aeditio. ¶ Et lucem habi,tat iaccessibilem. ¶ Quia vetus ĩterpres participium mutauit in verbum: copulandi particulam iiis adiecit q̃ sine ea particula Paulus profert. Paulus. Qῶσ ὀικῶμ ἀττρόσιτο p. ¶ Vulgata ,aeditio. Sed nec videri potest. ¶ Actiue potius videre. Paulus. ουAὲ ἐλεῖμ λύμαται» ¶ Vulgata aeditio. Praecipe non sublime sapere neque sperare in icerto diuitiarum et Dictum em alias nostros malle dici ac iterpretari non alte sentireq̃ sapere. Et in abstracto dicitur i icertitudine. Paulus. παράγγε λλε μὴ ὐχηλοφροpεῖφ, μηλὲ ηλπηχέμαι ἐττὶ » ωλούτου αAuλότuTι. ¶ Vulgata aeditio. Qui praestat nobis omnia abunde ad fruendum. ¶ Qui abhorrent dicere nos frui rebus a diuina etiam prouidentia elargitis: minus anxie fe rent si dicatur ad delectabiliter vtendumi siue delectabilem vsum. Et certe etiam spirituales virit bona dei participantes: fruuntur at non tam rebus q̃ deo largitore rerum. Cum enim cogito deum tam suauiter omnia dispensare: quomodo in re domini mei et praesertim quia domini mei tam be nigne tam suauiter omni rei prouidentis possim non delectari? Et hoc e: in hoc mundo frui. et ad sic fruendum: omnia nobis suppeditat deus. Alius autem est in hoc saeculo fruendi modus: de quo hic non est sermo. Et ad hunc et ad illum fruendi modum: vocabulum hoci plane em homony mum. Paulus. τῷ τταρέ χομτι Ἀμίμ πάυτα ωλουσίωσ, ἐἰς ἀτόλαυGιh. ¶¶Vulgata aeditio. Communicarei thesaurizare sibi fundamentum bonum in futurum vt apprehendant veram vitam. (¶In principio Paulus nominibus et parti cipiis vtituri non verbis. Et codices sunt qi fi nelaeternam habent vitam non veram. Paulus. κοιμωρικοὺς, ἀττοθησαυρίφορτας ἐαυrοῖσ oεnέλιομ Καλὸν ὲις τὸ μέλλομ, ἰμα ἐπιλάθωυται τησ ἀιωυίου τou͂ς¶ Vulgata aeditio. Deuitans prophanas vocum nouitates. ¶ Dici potest: euertens. Et vocabulum quo vtitur Paulus: potius vanas inanesque voces significat q̃ nouas voces aut vocum nouitates. Paulus. ἐκτρετόμεμοσ τασ pεφή λουσ Rεμοφωμίας. ει Vulgata aeditio. Et oppositiones falsi nominis scientiae. ( Quod hic vetus interpres dicit seientiae: gnosis a Paulo dicitur vnde Gnostici dicunturi quaẽ perperam et falso gnosis nominatur. Cum gnosi et verae quidem gnosi:sit opposita. Paulus. καὶ ααμτιθέσεισ τη̃σ Peυλωρύμον γρώσςωσ. Et apud Grecos subscriptio habetur. Prima ad Timotheum  scripta Laodociae.  \pend\pstart CCOMMENTARIORVM IN EPISTOLAS BEATISSIME PAV. APOSTOLIJ DECIMO LIBRO QVI EST EPISTOLAE PRIMAE AD TIMOTHEVM EXPLETO: CHRISTO DOMINO CVM PATRE ET SPIRITVSANCTOIVNI DEO OMNIVM LARGITORI BONORVMI GRATIAEIN PERPETVVM AMEN.  \pend
\endnumbering
\end{pages}
\end{document}
        